% Modernised reading — fonds Alexandre Grothendieck, shelfmark 80
% (pages 1-72, the whole folder).
%
% This is an INTERPRETATION, not a transcription. It re-reads the pages in
% today's notation and vocabulary, states the hypotheses the manuscript leaves
% implicit, and drops the critical apparatus so that the mathematics reads
% straight through. Everything the brackets and underlines carried moves into a
% footnote. It is held to being mathematically correct as it stands; where the
% manuscript is loose or mistaken, the true statement is what appears here and
% the footnote says what the page has. In any dispute of substance the
% transcription governs: whoever cites Grothendieck must cite batch-01.fr.tex
% (pp. 1-20), batch-02.fr.tex (pp. 21-40), batch-03.fr.tex (pp. 41-60) or
% batch-04.fr.tex (pp. 61-72).
%
% Pass: Opus 5.5 (claude-opus-5-5), 2026-10-10 — first pass, unchecked against the pages by a human.
% The folder was read whole, its four batches together; this is one document
% for the shelfmark. It was made on Opus 5.5 and has not been measured against
% the reference reading of folder 115 (made on Opus 5). The literature named in
% the footnotes is cited from memory and was not checked against the sources.
% Folders 81, 88 and 89 of the same group, read in their modernised readings,
% served as models; cross-references to them are ours and said so.
%
% Scope. Pages 1, 8, 15, 19, 30, 55, 59, 71 are blank sheets. Pages 21, 23,
% 25, 29 are photocopies of a letter in German unrelated to the mathematics,
% and page 27 a typed administrative notice dated 13 May 1981; the
% transcription does not reproduce them, nor does this reading (RIGHTS.md).
% Pages 9, 11-14, 56-58 are drawings, given only as the transcription
% describes them.
%
% Spine. Arrangements of pseudo-lines in the real projective plane, seen as
% maps (« cartes »), and the question how much of the map is determined by
% its 1-skeleton with the decomposition of the edges into pseudo-lines.
%   pp. 2-7     counting arrangements with a distinguished line: cut along it,
%               a disc with a stable system of chords; orbit-stabiliser counts.
%   pp. 9, 11-14 drawings (a nine-curve arrangement; polyhedra; a rectilinear
%               construction with inequalities).
%   p. 10       unrelated: Stein factorisation and base change.
%   pp. 16-18   real hyperplane arrangements: sign vectors, closures of strata.
%   pp. 20-28   the « crucial » arrangements: bigonal lines, the types I_N,
%               I_{p,q}, II'_{p,q}, II_4, II_5, III_6, IV_7; a Théorème (p. 24),
%               the lattice case and a Proposition (p. 26), a Lemme (p. 28).
%   pp. 31-38   three pseudo-lines: cube, octahedron, their « gauche »
%               (projective) quotients, orientations via contracted products,
%               the equivalence (16) and the cube of functors (17).
%   pp. 39-54   oriented 1-dimensional arrangements; the « Th. » (p. 39, the
%               functor (18) is fully faithful), its proof by reconstruction
%               (pp. 41-49, conditions (A1)-(A4)), the canonical orientation of
%               triangles (pp. 49-54, (A5), (A6)), the « Théorème » of p. 54.
%   pp. 56-58   drawings with region counts.
%   pp. 60-66   second start, « Système de Pseudodroites »: definition,
%               orientation cover, the bijections phi_{s;D,D'}, the local
%               system Omega(s).
%   pp. 67-70   third start with oriented pseudo-lines: rho_{D,s}, axioms
%               a)-f), « Th » (p. 70).
%   p. 72       table of small arrangements, projective vs combinatorial.
% The three reconstructions (pp. 41-54, 60-66, 67-70) are kept distinct.
%
% Notation fixed for the document.
%  - X the surface (a real projective plane), 𝒳 the map; (D_i)_{i in I} its
%    pseudo-lines, I the set of ALL lines (as pp. 41-54); a triple of lines is
%    J in 𝔓_3(I). Page 39 swaps the letters (J all lines, I a triple);
%    footnoted. In pp. 31-38, which treat three lines only, I is the
%    three-element index set, as on the page.
%  - In pp. 2-7 the page's own (X, Σ, Δ): Σ = (Δ_i)_{i in I} the family of
%    lines, Δ = Δ_{i_0} the distinguished one; card I = n+1, I_0 = I \ {i_0}.
%  - « stable » (index st) = simple: every vertex on exactly two lines.
%  - (K, (D_i)) a 1-dimensional arrangement (axioms of p. 40); K_J the
%    sub-arrangement of the lines in J; nu_s the number of lines through s
%    (the vertex has order 2 nu_s).
%  - A two-element set is a {±1}-torsor; ∧ the contracted product, 1 the
%    trivial torsor; or(I) the two orientations of a finite set I (for card
%    I = 3, its two circular orders: folder 89's omega_I).
%  - A polygonal structure on a finite set E (card >= 3): a pair {u, u^{-1}}
%    of inverse circular permutations (folder 81's bi-circulation); 𝒞_E the
%    set of circular permutations of E.
%  - 𝔷_s (the page's 𝔷^1_s with a doubtful superscript, dropped) the
%    polygonal structure on the lines through s; vec-𝔷_s the pair of inverse
%    circular orders on the oriented lines through s.
%  - The doubtful superscript on xi (pp. 5-7) is dropped.
%  Collisions not corrected, the sections not mixing: S is the pinning polygon
%  (pp. 3-7), the sphere (pp. 16-18, bold) and the vertex set of K (p. 41);
%  Omega_i are orientations of a bigon (pp. 34-37), Omega(D), Omega(s)
%  orientations of a pseudo-line and at a vertex (pp. 61-70).
%
% Substantive departures, each footnoted where it occurs:
%  p. 2    the counts « n = 5: C(n) has one element, fibre of card 6 » hold
%          for simple (stable) arrangements only; surjectivity of beta'
%          (asked with « ? ») is Levi's enlargement lemma (ours).
%  p. 5    the induction formula, « sans doute » on the page, is a general
%          fact (an equivariant map to G/H is induced from its fibre).
%  p. 6    Ex. 1: the orbit of order 2 is under D_S (order 4n), not D_{I_0}.
%  p. 9    region counts: 3·7+4·24+5·6 = 147 is odd, so they cannot all be
%          faces of a map counted with their sides; flagged, not resolved.
%  p. 10   the Stein factorisation is the FINAL, not initial, object of the
%          category of Stein arrows over Z -> X, with the arrows as drawn.
%  p. 16   V^* is V \ {0}; D_H ≃ F_3 only after a choice of coorientation.
%  p. 18   « Σ est une partie ouverte de E » is false in general (the page's
%          own struck example, three planes through a line); true exactly
%          when, at each x != 0, the forms vanishing at x are independent;
%          the Inf statement follows only then.
%  p. 22   I_N has N vertices, not the margin's N+1.
%  p. 24   I'_{3,2} (II_{3,2} pp. 26-28), undefined, read as II'_{3,2};
%          then IV_7 contains a II'_{3,2} (explicit), so the parenthesis
%          « s'il ne contient ni I'_{3,2} ni II_5, c'est III_6 ou IV_7 » must
%          read « II_4 ou III_6 »; the boxed statement is unaffected. Same for
%          the NB of p. 20.
%  p. 26   « pas de type II_4 » read as « K n'est pas de type II_4 » (II_5
%          contains II_4 and is said to be lattice-like on p. 24).
%  p. 33   C -> C~, O -> O~ are full (onto on automorphisms) and NOT faithful
%          (kernel {1, a}); the page says the reverse, with an argument that
%          proves non-faithfulness.
%  p. 38   in the cube (17), if it commutes, gr.oct -> gr.oct.gauches cannot
%          be an equivalence: it kills the antipody and its image has index 2.
%  p. 40   « octogonal » read octaédral.
%  p. 42   the margin's a priori c_2 <= c(K): true because X is connected and
%          not a sphere (argument ours).
%  p. 49   (A5) names two different conditions (pp. 49 and 51).
%  p. 50   « concourantes » read « non concourantes ».
%  p. 60   the definition needs each class of edges to pass at most once
%          through each vertex (a simple closed curve); « n <= 6 »: all
%          arrangements of at most 8 pseudo-lines are stretchable
%          (Goodman-Pollack 1980), and nine is where non-stretchable ones
%          start (Ringel 1956).
%  p. 64   a) « simplement transitif » read « système transitif ».
%  p. 72   4 lines in general position: the projective stabiliser of the SET
%          of lines is 𝔖_4 too; « Id » holds for the transformations fixing
%          each line. 5 lines in general position are NOT projectively
%          standard (2 moduli); the « pas » the page strikes is needed.
%
% Links the pages do not write (ours, and said so in the body): the chord
% systems of pp. 3-7 and the drawings of folder 81 (pp. 9-11); the labels of
% p. 9 recur on p. 13; the octahedron of p. 11 (three great circles) is the
% octahedron of p. 31 and 48 = 2·4! its symmetry group, which p. 38's margin
% writes 𝔖_4 × {±1}; the six-line configurations of pp. 20, 24, 53, 63-64
% are one and the same, the complete quadrangle III_6; u_{s,Δ} of p. 45 is
% the rho_{D,s} of p. 67; the axioms of p. 67 restate (30), (A3), (A4), (22).
%
% Announced and never established: the « Théorème » of p. 24 and the
% Proposition (lattice case) and Proposition of p. 26; the Lemme of p. 28
% (an « Idée » only); the « Th. » of p. 39, whose proof he says he has not
% checked in the case of n concurrent lines and one transversal (p. 49);
% (A1)-(A6) in the 2-dimensional case (« bien entendu »); the reduction of
% (A1)-(A6) and the automaticity of (23) (p. 54); the end of the condition on
% r (p. 63); item c) of p. 64; the uniqueness asserted at the top of p. 65;
% the « Th » of p. 70; the realisation for n <= 6 (p. 60, « on montre »);
% the base-change question of p. 10.
\documentclass[11pt,a4paper]{article}
\input{../preamble/grothendieck}

\folder{80}
\pages{1}{72}
\dating{[à partir de 1981]}
\watermark{Édition de démonstration\\interprétation personnelle de l'œuvre}
\foldertitle{Systèmes de pseudo-droites : notes manuscrites (s.d.) — lecture modernisée du dossier entier}

\begin{document}

\section*{Résumé}

\begin{resume}
Tracez sur une feuille des courbes qui ne soient pas forcément droites, mais
qui se comportent comme des droites : chacune coupe chacune des autres en un
point et un seul, et aucune ne se ferme sur elle-même comme un cercle. Pour
que la règle « un point et un seul » vaille sans exception, même pour deux
courbes qui sembleraient parallèles, on travaille sur le plan projectif réel
--- un disque dont on recolle chaque point du bord à son opposé ---, où deux
droites ordinaires se coupent toujours. Ces courbes s'appellent des
\emph{pseudo-droites}, et le dessin qu'elles forment, avec ses points de
croisement, ses segments et ses régions, est une \emph{carte} au sens de
Grothendieck : un réseau de sommets, d'arêtes et de faces tracé sur une
surface.

La question qui court à travers ces soixante-douze pages est celle-ci. Que
faut-il savoir pour reconstituer un tel dessin ? Imaginez une ville dont on
ne connaît que la liste des rues et, pour chaque rue, l'ordre des carrefours
qu'on rencontre en la parcourant. Peut-on en redessiner le plan, avec ses
pâtés de maisons ? Grothendieck appelle cette donnée un arrangement
« de dimension $1$ » --- le squelette fait d'arêtes, découpé en
pseudo-droites --- et le dessin complet un arrangement « de dimension $2$ ».
Sa réponse, qu'il énonce sans l'avoir entièrement démontrée, est :
presque toujours, oui. Il y a un obstacle, qu'il isole avec
soin : une rue qui n'a que deux carrefours (un « bigone ») laisse une
ambiguïté que rien ne lève, comme deux chemins parallèles entre les mêmes
places. Hors de ce cas, le squelette suffit, et même les symétries du
squelette sont toutes des symétries du dessin.

Pour y parvenir, il passe par le cas de trois pseudo-droites --- qui découpent
le plan projectif en quatre triangles, la moitié d'un octaèdre ---, et par une
comptabilité minutieuse des \emph{orientations} : un objet à deux états,
comme les deux sens de parcours d'une boucle, et la façon dont ces paires
d'états se combinent. Il cherche ensuite à lire, en chaque carrefour, l'ordre
circulaire dans lequel les rues en partent, à partir des seules rues qui n'y
passent pas : une rue lointaine, coupée successivement par toutes celles du
carrefour, en donne l'ordre, comme l'horizon donne l'ordre des chemins qui
partent d'un col. Il réécrit cette reconstruction trois fois, de trois points
de départ différents.

Le dossier contient aussi un dénombrement des arrangements munis d'une droite
distinguée, une étude des arrangements d'hyperplans de l'espace réel
décrits par les signes que prend un point par rapport à chaque hyperplan, une
liste des plus petites configurations qui ne sont pas « dégénérées », et
un tableau qui compare, pour un, deux, \ldots, six droites, ce que voit la
géométrie projective et ce que voit la combinatoire. Une page isolée, sans
rapport avec le reste, porte sur la factorisation de Stein.

On y voit une manière de travailler. Grothendieck signale lui-même ce qu'il
n'a pas vérifié --- « Je n'ai donc pas vérifié que la proposition énoncée
soit correcte » --- et ce qu'il ne fait que soupçonner --- « je soupçonne
que cette dernière est automatique ». Il écrit aussi « j'ai maintenant envie
de montrer », et la suite du texte suit cette envie plutôt qu'un plan. Et
quand une approche s'embarrasse de cas particuliers, il ne la raccommode
pas : il reprend tout, d'un autre point de départ.

Les noms sous lesquels chercher la suite : arrangements de pseudo-droites,
lemme d'extension de Levi, étirabilité ; matroïdes orientés, vecteurs de
signes, théorème de représentation topologique de Folkman et Lawrence ;
quadrangle et quadrilatère complets ; systèmes de rotation des plongements
de graphes.

\keywords{pseudoline arrangement, real projective plane, Levi enlargement lemma, orbit-stabilizer counting, stretchability, Stein factorization, monotone-light factorization, oriented matroid, covector, face poset, Folkman-Lawrence representation theorem, simple arrangement, complete quadrangle, complete quadrilateral, hemi-octahedron, contracted product of torsors, rotation system, orientation double cover, local system}
\end{resume}

\pagerange{1}{72}
\section*{Le fil du dossier, et les conventions}

\subsection*{Les stations}

Le dossier n'est pas daté de sa main ; l'inventaire le date « à partir de
1981 », ce qu'appuie la notice administrative dactylographiée de la page 27,
datée du 13 mai 1981, et le range dans le groupe « Géométrie et topologie
combinatoire ». Sa pagination propre court sur deux séries : $1$ à $6$ aux
pages 2 à 7, et $1$ à $12$ sur les rectos des pages 31 à 53, les versos
continuant le texte. Ses équations numérotées (1) à (34) et ses conditions
(A1) à (A6) courent des pages 31 à 54.

\begin{enumerate}
  \item \emph{Pages 2 à 7}, « Classifications pseudo-droites » : les
  arrangements munis d'une pseudo-droite distinguée, découpés le long de
  celle-ci en un disque muni d'un système de cordes ; dénombrements par
  orbites et stabilisateurs.
  \item \emph{Pages 9 et 11 à 14}, des dessins : un arrangement de neuf
  courbes, des polyèdres, une construction rectiligne avec inégalités. La
  page 10, au milieu, est étrangère au sujet : la factorisation de Stein.
  \item \emph{Pages 16 à 18}, « Système d'hyperplans (projectifs) réels » :
  vecteurs de signes, adhérence des strates, ordre sur les signes.
  \item \emph{Pages 20 à 28}, les arrangements « cruciaux » : pseudo-droites
  bigonales, une liste de types, un théorème (page 24), le cas
  « latticiel » (page 26), un lemme d'ordonnancement (page 28).
  \item \emph{Pages 31 à 38}, « Arrangements de 3 pseudo-droites » : cube et
  octaèdre combinatoires, leurs quotients « gauches », les orientations, le
  cube de foncteurs (17).
  \item \emph{Pages 39 à 54}, arrangements de dimension $1$ orientés, le
  théorème de pleine fidélité de la page 39, sa démonstration par
  reconstruction (pages 41 à 49), l'orientation canonique des triangles
  (pages 49 à 54) et le théorème de la page 54.
  \item \emph{Pages 56 à 58}, des dessins avec décompte des régions.
  \item \emph{Pages 60 à 66}, « Système de Pseudodroites » : un second
  départ, par les ensembles d'orientations $\Omega(D)$ et un système local.
  \item \emph{Pages 67 à 70}, un troisième départ, par les pseudo-droites
  orientées et les permutations circulaires $\rho_{\vec{D},s}$.
  \item \emph{Page 72}, un tableau des arrangements de une à six droites.
\end{enumerate}

Les pages 1, 8, 15, 19, 30, 55, 59 et 71 sont blanches. Les pages 21, 23, 25
et 29 sont des photocopies d'une lettre en allemand étrangère aux
mathématiques, la page 27 la notice déjà dite ; ni la transcription ni cette
lecture ne les reproduisent.

\subsection*{Trois reconstructions, gardées distinctes}

La même question --- retrouver la carte à partir de son squelette découpé en
pseudo-droites --- est attaquée trois fois, et les trois rédactions ne se
citent pas.
\begin{itemize}
  \item \emph{Pages 39 à 54} : on se donne en outre une \emph{orientation}
  $\omega$ de chaque triple de pseudo-droites non concourantes ; on
  reconstruit en chaque sommet la paire d'ordres circulaires
  $\vec{\mathfrak{Z}}_s$ et, le long de chaque arête, un transport
  $\psi_{\vec{a}}$ ; puis on montre que $\omega$ est elle-même déterminée par
  le squelette.
  \item \emph{Pages 60 à 66} : on part des ensembles $\Omega(D)$ des deux
  orientations de chaque pseudo-droite et de bijections
  $\varphi_{s;D,D'} : \Omega(D) \simeq \Omega(D')$ attachées aux sommets, et
  on en tire un système local $\Omega(s)$ sur le squelette.
  \item \emph{Pages 67 à 70} : on part des pseudo-droites orientées et d'une
  permutation circulaire $\rho_{\vec{D},s}$ des pseudo-droites orientées
  issues de $s$, pour chaque $\vec{D}$ ne passant pas par $s$, soumise à des
  axiomes a) à f).
\end{itemize}
Ce sont trois habillages de la même construction --- la permutation
$\vec{u}_{s,\vec{\Delta}}$ de la page 45 est exactement la
$\rho_{\vec{\Delta},s}$ de la page 67 ---, mais chacun pose ses propres
données et ses propres axiomes, et on les lit séparément.

\subsection*{Conventions}

\begin{itemize}
  \item \emph{Arrangement.} Un arrangement de pseudo-droites, ou
  \emph{système de pseudo-droites}, est une carte $\mathcal{X}$ sur une
  surface $X$ homéomorphe au plan projectif réel, dont les arêtes se
  répartissent en pseudo-droites (définition précise : page 60). Ses
  pseudo-droites sont notées $(D_i)_{i \in I}$, et $I$ est toujours
  l'ensemble de \emph{toutes} les pseudo-droites ; un triple est une partie
  $J \in \mathfrak{P}_3(I)$\footnote{La page 39 fait l'inverse : $J$ y
  indexe toutes les pseudo-droites et $I \in \mathfrak{P}_3(J)$ un triple.
  Les pages 41 à 54 adoptent la convention suivie ici. Aux pages 31 à 38,
  qui ne traitent que de trois pseudo-droites, $I$ est l'ensemble à trois
  éléments qui les indexe, comme sur la page.}. Aux pages 2 à 7 on garde la
  notation de la page, $(X, \Sigma, \Delta)$ : $\Sigma = (\Delta_i)_{i \in I}$
  la famille des pseudo-droites, $\Delta = \Delta_{i_0}$ la pseudo-droite
  distinguée.
  \item \emph{Simple.} L'indice « $\mathrm{st}$ » (« stable ») des pages 2 à
  7 désigne les arrangements \emph{simples} : par chaque sommet passent
  exactement deux pseudo-droites. C'est le nom actuel, et c'est bien la
  stabilité par petite déformation\footnote{La page ne définit pas
  « stable » ; l'identification est nôtre, et c'est la seule qui rende
  justes les nombres des pages 2 à 7.}.
  \item \emph{Bigone, monogone.} Une pseudo-droite qui ne porte que deux
  sommets est \emph{bigonale} : sur le squelette, c'est un polygone à deux
  côtés. Une pseudo-droite qui n'en porte qu'un est un monogone, ce qui
  n'arrive que si toutes les pseudo-droites sont concourantes.
  \item \emph{Arrangement de dimension $1$.} C'est le squelette : un
  $1$-complexe $K$ et une famille $(D_i)_{i \in I}$ de sous-complexes
  polygonaux, avec les axiomes de la page 40. Pour $J \subset I$, $K_J$ est le
  sous-arrangement formé des $D_i$, $i \in J$, où l'on oublie les sommets
  qui ne sont pas des intersections de ces $D_i$. On note $\nu_s$ le nombre
  de pseudo-droites passant par un sommet $s$ ; l'ordre du sommet est
  $2\nu_s$.
  \item \emph{Torseurs.} Un ensemble à deux éléments est un torseur sous
  $\{\pm 1\}$ ; $\wedge$ est le produit contracté, $\underline{1}$ le
  torseur trivial, et deux copies d'un même torseur ont un produit
  canoniquement trivial. Pour un ensemble fini $I$, $\mathrm{or}(I)$ est
  l'ensemble de ses deux orientations (ordres totaux à permutation paire
  près) ; pour $\operatorname{card} I = 3$, ce sont ses deux ordres
  circulaires --- l'ensemble que le dossier 89 note $\omega_I$. On garde le
  signe $\wedge$ des pages, qui est celui du dossier 89.
  \item \emph{Polygones.} Une \emph{structure polygonale} sur un ensemble
  fini $E$ de cardinal $\geqslant 3$ est une paire $\{u, u^{-1}\}$ de
  permutations circulaires inverses l'une de l'autre --- la
  « bi-circulation » du dossier 81 (page 41), et la troisième description
  d'un polygone combinatoire du dossier 89 (pages 19 et 20). $\mathcal{C}_E$
  est l'ensemble des permutations circulaires de $E$. Les sommets portés par
  une pseudo-droite forment un polygone, donc reçoivent une telle
  structure.
\end{itemize}

Restent des collisions de lettres qu'on ne corrige pas, les sections ne se
mêlant pas : $S$ est le polygone d'épinglage (pages 3 à 7), la sphère
(pages 16 à 18, en gras) et l'ensemble des sommets de $K$ (page 41) ;
$\Omega_i$ est l'ensemble des orientations d'un bigone (pages 34 à 37),
$\Omega(D)$ et $\Omega(s)$ ceux d'une pseudo-droite et d'un sommet (pages 61
à 70).

\subsection*{Ce que seule une lecture d'ensemble peut dire}

Des liens qu'aucun feuillet n'écrit, et qui sont de nous. La configuration
de six pseudo-droites joignant deux à deux quatre points --- le
\emph{quadrangle complet} --- revient quatre fois : dessinée page 20 (« $6$
en position particulière »), nommée $\mathrm{III}_6$ page 24, rencontrée
comme cas résiduel page 53, et pages 63 et 64 comme celui où l'on tombe
« par malheur ». L'octaèdre tracé page 11 comme une sphère coupée de trois grands
cercles est l'octaèdre combinatoire de la page 31, et le nombre
$48 = 2 \cdot 4!$ qui l'accompagne est l'ordre de son groupe de symétries,
$\mathfrak{S}_4 \times \{\pm 1\}$, que la marge de la page 38 écrit. Les
étiquettes des points de la page 9 se retrouvent sur la construction de la
page 13. Les axiomes de la page 67 réécrivent, pour les $\rho_{\vec{D},s}$,
les relations (30), (A3), (A4) et (22) des pages 42 à 47. Enfin, les disques
coupés de cordes des pages 3 à 7 sont de la même famille que ceux que
dessine le dossier 81 (pages 9 à 11), mais ici les cordes se croisent toutes
deux à deux : ce sont les pseudo-droites elles-mêmes, vues après découpage.

\subsection*{Ce que le dossier annonce sans l'établir}

Le théorème de la page 24, la proposition sur le cas latticiel et la
proposition sur l'adjonction d'une pseudo-droite (page 26), le lemme de la
page 28, dont seule une « Idée » est écrite. Le « Th. » de la page 39, dont
Grothendieck dit lui-même (page 49) ne pas avoir vérifié la démonstration
dans le cas de $n$ pseudo-droites concourantes coupées par une seule autre ;
les conditions (A1) à (A6), affirmées vraies pour un arrangement de
dimension $2$ (« bien entendu ») ; la simplification de ces conditions et le
caractère automatique de la relation d'Euler-Poincaré (page 54). La fin de
la condition sur le point $r$ (page 63), l'item c) de la page 64,
l'unicité qu'affirme le haut de la page 65, le « Th » de la page 70. La
réalisation par de vraies droites pour $n \leqslant 6$ (page 60, « on
montre »). La question du changement de base de la page 10.

\pagerange{2}{7}
\section*{I. Arrangements à pseudo-droite distinguée (pages 2 à 7)}

\pagerange{2}{2}
\subsection*{Les ensembles de classes (page 2)}

Soit $I$ un ensemble à $n+1$ éléments, $i_0 \in I$ et $I_0 = I \smallsetminus
\{i_0\}$. On note $\mathrm{C}(I)$ l'ensemble des classes d'isomorphisme
d'arrangements de pseudo-droites $(X, \Sigma = (\Delta_i)_{i \in I})$ indexés
par $I$, et $\mathrm{C}_{\mathrm{aff}}(I_0)$ celui des arrangements
\emph{affines} indexés par $I_0$, c'est-à-dire tracés dans le complémentaire
d'une pseudo-droite $\Delta$ qu'on regarde comme la droite à
l'infini\footnote{L'indice « $\mathrm{aff}$ » est une lecture incertaine :
il l'écrit comme un petit a suivi de deux hampes doublées. La lecture
« affine » est la seule qui convienne à l'énoncé
$\mathrm{C}(I) \simeq \mathrm{C}_{\mathrm{aff}}(I_0)$.}. Choisir pour droite
à l'infini la pseudo-droite d'indice $i_0$ donne une bijection
\[
\mathrm{C}(I) \simeq \mathrm{C}_{\mathrm{aff}}(I_0).
\]
Le groupe symétrique $\mathfrak{S}_I$ agit en changeant les indices, et l'on
pose $\mathrm{C}(n+1) = \mathrm{C}(I)/\mathfrak{S}_I$ (arrangements de $n+1$
pseudo-droites, sans indices) et $\mathrm{C}_{\mathrm{aff}}(n) =
\mathrm{C}_{\mathrm{aff}}(I_0)/\mathfrak{S}_{I_0}$ ; ce dernier est
l'ensemble des classes de triples $(X, \Sigma, \Delta)$ avec $\Delta \in
\Sigma$, la pseudo-droite distinguée seule gardant son nom. On a deux
applications :

\begin{tikzcd}
& \mathrm{C}_{\mathrm{aff}}(n) \arrow[dl, "\beta"'] \arrow[dr, "\beta'"] & \\
\mathrm{C}(n+1) & & \mathrm{C}(n)
\end{tikzcd}

$\beta$ oublie quelle pseudo-droite est distinguée, $\beta'$ efface
$\Delta$\footnote{Les étiquettes de la page sont « oubli du $\Delta$ » pour
$\beta$, en partie illisible, et « gommage de $\Delta$ », écrit au-dessus de
« oubli de $\Delta$ » biffé, pour $\beta'$. La page porte aussi une flèche
$\alpha$ de $\mathrm{C}(I)$ vers $\mathrm{C}_{\mathrm{aff}}(n)$, « oubli du
$I$-épinglage sauf $\Delta_{i_0} = \Delta$ », qui est le passage au quotient
par $\mathfrak{S}_{I_0}$. Un premier état du diagramme est biffé.}.

$\beta$ est surjective, et sa fibre en la classe de $(X, \Sigma)$ est
l'ensemble des orbites de pseudo-droites sous le groupe d'automorphismes
$G_{X,\Sigma}$ : $\Sigma / G_{X,\Sigma}$. $\beta'$ est surjective aussi ---
la page le demande avec un point d'interrogation --- : tout arrangement se
prolonge par une pseudo-droite supplémentaire, et même par une pseudo-droite
en position générale ; c'est le \emph{lemme d'extension de Levi}
(1926)\footnote{La réponse et le nom sont de nous ; la page n'en cite
aucun.}. Sa fibre est l'ensemble des « types de positions relatives » d'une
nouvelle pseudo-droite par rapport à l'arrangement donné.

Tout cela vaut aussi pour les arrangements simples : la partie
$\mathrm{C}_{\mathrm{aff,st}}(I_0) = \mathrm{C}_{\mathrm{st}}(I)$ de
$\mathrm{C}_{\mathrm{aff}}(I_0) = \mathrm{C}(I)$ est stable par
$\mathfrak{S}_I$, et en passant aux quotients on trouve
$\mathrm{C}_{\mathrm{aff,st}}(n)$, $\mathrm{C}_{\mathrm{st}}(n)$,
$\mathrm{C}_{\mathrm{st}}(n+1)$ et les mêmes diagrammes. Pour $n = 5$, la
page donne : $\mathrm{C}_{\mathrm{st}}(5)$ n'a qu'un élément et la fibre de
$\beta'$ en cet élément a $6$ éléments ; au-dessus de
$\mathrm{C}_{\mathrm{st}}(6)$, les fibres de $\beta$ ont $1$, $1$, $2$ et $2$
éléments\footnote{La page écrit ces nombres pour $\mathrm{C}(n)$ et
$\mathrm{C}(n+1)$, sans l'indice ; ils ne sont justes que pour les
arrangements simples --- il y a, par exemple, plusieurs arrangements non
simples de cinq pseudo-droites. Les deux comptes donnent le même total,
$6 = 1 + 1 + 2 + 2$, qui est le cardinal de
$\mathrm{C}_{\mathrm{aff,st}}(5)$ écrit page 7 : c'est une vérification de
cohérence que nous faisons, non une vérification des nombres eux-mêmes.}.

\pagerange{3}{4}
\subsection*{Couper le long de la droite distinguée (pages 3 et 4)}

Soit $\xi = \mathrm{cl}(X, \Sigma) \in \mathrm{C}_{\mathrm{st}}(I) \simeq
\mathrm{C}_{\mathrm{aff,st}}(I_0)$, avec $\operatorname{card} I_0 \geqslant 3$.
Les pseudo-droites $\Delta_i$, $i \in I_0$, coupent $\Delta$ en des points
distincts (l'arrangement est simple), rangés sur $\Delta$ dans un ordre
circulaire : $I_0$ hérite d'une \emph{structure polygonale} $\pi(\xi)$, qui
dépend de $\xi$.

Découpons $X$ le long de $\Delta$. On obtient un disque fermé $\mathbf{D} =
\mathrm{Déc}(X, \Delta)$, dont le bord $\partial \mathbf{D}$ est un revêtement
à deux feuillets de $\Delta$ ; chaque $\Delta_i$, $i \in I_0$, devient une
\emph{corde} de $\mathbf{D}$, dont les deux extrémités sont les deux points
de $\partial\mathbf{D}$ au-dessus de $\Delta_i \cap \Delta$ --- donc deux
points antipodaux du bord. On obtient un système de $n$ cordes deux à deux
sécantes, simple, indexé par $I_0$, et l'ensemble
\[
\widetilde{I}_0(\xi) = \{\text{extrémités des cordes sur } \partial\mathbf{D}\}
\]
est un polygone à $2n$ sommets, revêtement à deux feuillets du polygone
$I_0(\xi)$, l'antipodie échangeant les deux feuillets\footnote{La page
appelle ce disque $\mathbf{D} = \widetilde{X}$ et l'abréviation
« $\mathrm{Déc}$ » est la sienne. Elle note en marge « deux feuillets de
$\Delta$ ». La description des cordes comme joignant des points antipodaux
est explicitée par nous ; elle découle de ce que $\partial\mathbf{D} \to
\Delta$ est le revêtement double.}.

Fixons maintenant, une fois pour toutes, un polygone combinatoire $S$
d'ordre $2n$, et notons $a_S$ son antipodie (la rotation d'un demi-tour) et
$I_0 = S/\{1, a_S\}$, muni de la structure polygonale quotient $\pi_0$. Soit
$\widetilde{\mathrm{C}}_{\mathrm{st}}(S)$ l'ensemble des classes de disques
munis d'un système simple de cordes \emph{épinglé} par $S$, c'est-à-dire
dont l'ensemble des extrémités est identifié à $S$ par un isomorphisme de
polygones. L'application naturelle
\[
\widetilde{\mathrm{C}}_{\mathrm{st}}(S) \longrightarrow
\mathrm{C}^{!}_{\mathrm{aff,st}}(I_0) := \{\xi \in
\mathrm{C}_{\mathrm{aff,st}}(I_0) \mid \pi(\xi) = \pi_0\}
\]
a pour image la partie notée $\mathrm{C}^{!}$, formée des $\xi$ dont la
structure polygonale induite sur $I_0$ est $\pi_0$. Sa fibre en $\xi$
s'identifie à l'ensemble des isomorphismes de polygones $S \simeq
\widetilde{I}_0(\xi)$ au-dessus de $I_0$ ; elle a \emph{deux} éléments, et
on passe de l'un à l'autre par l'antipodie\footnote{Pour que la fibre ait
exactement deux éléments, il faut qu'un automorphisme de $(X, \Sigma)$ qui
fixe chaque pseudo-droite soit l'identité ; c'est le cas pour un arrangement
simple d'au moins quatre pseudo-droites, puisqu'il fixe alors chaque sommet,
intersection de deux pseudo-droites. Cette justification est de nous. La
page commence par écrire « dont la fibre en $\xi$ est un \ldots », biffé, et
le second argument de $\mathrm{C}_{\mathrm{aff,st}}(I_0, \ldots)$ est
illisible ; on lit la notation $\mathrm{C}^{!}$ d'après la page 4, qui
l'explique.}. Autrement dit,
\[
\widetilde{\mathrm{C}}_{\mathrm{st}}(S)/\{1, a_S\} \simeq
\mathrm{C}^{!}_{\mathrm{aff,st}}(I_0).
\]

\pagerange{5}{7}
\subsection*{Compter (pages 5 à 7)}

Le groupe diédral $\mathbb{D}_{I_0} \subset \mathfrak{S}_{I_0}$, groupe
d'automorphismes de la structure polygonale $\pi_0$, d'ordre $2n$, agit sur
$\mathrm{C}^{!}_{\mathrm{aff,st}}(I_0)$, et l'on a deux applications
\[
\alpha : \mathrm{C}_{\mathrm{aff,st}}(I_0) \longrightarrow
\mathrm{C}_{\mathrm{aff,st}}(n), \qquad
\alpha^{!} : \mathrm{C}^{!}_{\mathrm{aff,st}}(I_0) \longrightarrow
\mathrm{C}_{\mathrm{aff,st}}(n) \simeq
\mathrm{C}^{!}_{\mathrm{aff,st}}(I_0)/\mathbb{D}_{I_0}.
\]
Soit $\xi \in \mathrm{C}_{\mathrm{aff,st}}(n)$\footnote{La page note cet
élément $\xi$ surmonté d'un petit signe en exposant, lu « $q$ » ou « $4$ »
sans certitude ; on l'omet. Elle écrit $\mathrm{C}_{\mathrm{aff}}(n)$ là où,
les arrangements étant simples, il faut $\mathrm{C}_{\mathrm{aff,st}}(n)$.}.
La fibre $\alpha^{-1}(\xi)$ est une orbite de $\mathfrak{S}_{I_0}$, la fibre
$(\alpha^{!})^{-1}(\xi)$ une orbite de $\mathbb{D}_{I_0}$, et
\[
\alpha^{-1}(\xi) \simeq \mathfrak{S}_{I_0} \times^{\mathbb{D}_{I_0}}
(\alpha^{!})^{-1}(\xi).
\]
Ce n'est pas une conjecture --- la page écrit « sans doute » --- mais un fait
général : l'application $\xi' \mapsto \pi(\xi')$ de $\alpha^{-1}(\xi)$ vers
l'ensemble des structures polygonales de $I_0$ est
$\mathfrak{S}_{I_0}$-équivariante, cet ensemble est homogène de stabilisateur
$\mathbb{D}_{I_0}$, et un ensemble équivariant au-dessus de $G/H$ est induit
de sa fibre en $H$\footnote{La justification est de nous. La page note le
produit contracté $\mathfrak{S}_{I_0} \wedge^{\mathbb{D}_{I_0}} (\cdot)$ ; on
écrit $\times^{\mathbb{D}_{I_0}}$, la notation usuelle pour l'induction, et
pour ne pas confondre avec le $\wedge$ des torseurs.}. D'où
\[
\operatorname{card} \alpha^{-1}(\xi) = \frac{n!}{2n} \cdot
\operatorname{card} (\alpha^{!})^{-1}(\xi), \qquad \frac{n!}{2n} =
\frac{(n-1)!}{2},
\]
le nombre de structures polygonales sur un ensemble à $n \geqslant 3$
éléments.

Le stabilisateur dans $\mathbb{D}_{I_0}$ d'un élément de
$(\alpha^{!})^{-1}(\xi)$ s'identifie au groupe $G_\xi = \mathrm{Aut}(X,
\Sigma, \Delta)$ des automorphismes de l'arrangement qui fixent $\Delta$ :
ils agissent fidèlement sur les indices, comme on l'a dit, et respectent
nécessairement l'ordre des points sur $\Delta$. Si $N(\xi) = \operatorname{card}
\mathrm{Aut}(X, \Sigma, \Delta)$, on a donc
\[
\operatorname{card} (\alpha^{!})^{-1}(\xi) = \frac{2n}{N(\xi)}, \qquad
\operatorname{card} \widetilde{\alpha}^{-1}(\xi) = \frac{4n}{N(\xi)},
\]
où $\widetilde{\alpha}$ est la composée $\widetilde{\mathrm{C}}_{\mathrm{st}}(S)
\to \mathrm{C}^{!}_{\mathrm{aff,st}}(I_0) \to \mathrm{C}_{\mathrm{aff,st}}(n)$,
de degré $2$ puis $\alpha^{!}$ ; c'est le quotient par le groupe diédral
$\mathbb{D}_S$ d'ordre $4n$ de $S$, et la formule est celle des orbites et
des stabilisateurs.

\textbf{Exemples.} Pour $n = 3$, $\mathrm{C}_{\mathrm{aff,st}}(3)$ a un
élément, $\mathrm{Aut}(X, \Sigma, \Delta) \simeq \mathfrak{S}_3$, et
$\widetilde{\mathrm{C}}_{\mathrm{st}}(S)$ est une seule orbite de
$\mathbb{D}_S$, à $12/6 = 2$ éléments\footnote{La page dit « une seule orbite
sous $\mathbb{D}_{I_0}$, d'ordre $12/6 = 2$ ». Le nombre $12 = 4n$ est
l'ordre de $\mathbb{D}_S$, et c'est sous $\mathbb{D}_S$ que
$\widetilde{\mathrm{C}}_{\mathrm{st}}(S)$ est une orbite ; sous
$\mathbb{D}_{I_0}$, d'ordre $6$, c'est $\mathrm{C}^{!}$ qui en est une, à un
seul élément.}. Pour $n = 4$, $\mathrm{C}_{\mathrm{aff,st}}(4)$ a un
élément, $\mathrm{Aut}(X, \Sigma, \Delta) \simeq \mathbf{Z}/2\mathbf{Z}$, et
$\operatorname{card} \widetilde{\mathrm{C}}_{\mathrm{st}}(S) = 16/2 = 8$. Pour
$n = 5$, $\operatorname{card} \mathrm{C}_{\mathrm{aff,st}}(5) = 6$, et la page
s'arrête là. Ces groupes sont les stabilisateurs d'une droite dans les
groupes d'automorphismes $\mathfrak{S}_4$ de quatre droites en position
générale et diédral d'ordre $10$ de cinq droites, que donne le tableau de
la page 72.

\textbf{Résumé de la page 7.} Grothendieck récapitule.
$\widetilde{\mathrm{C}}_{\mathrm{st}}(S)$ est l'ensemble des classes de
disques munis d'un système simple de cordes épinglé par le polygone $S$
d'ordre $2n$ ; il dépend fonctoriellement de $S$ pour les isomorphismes de
polygones, donc $\mathbb{D}_S$ y agit. Pour $\widetilde{\xi} \in
\widetilde{\mathrm{C}}_{\mathrm{st}}(S)$ correspondant à $(X, \Sigma,
\Delta)$, il demande si
\[
(\mathbb{D}_S)_{\widetilde{\xi}} \simeq \mathrm{Aut}_{\mathrm{ist}}(\mathbf{D},
\widetilde{\Sigma}) \simeq \mathrm{Aut}(X, \Sigma, \Delta),
\]
le premier terme étant le stabilisateur, le deuxième le groupe des
automorphismes à isotopie près du disque à cordes --- ce qu'on vient
d'utiliser\footnote{L'indice « $\mathrm{ist}$ » est lu ainsi, sans
certitude, et glosé en marge « (\ldots{} des cat.\ isotopiques) ». La page
fait précéder la formule d'un « alors ? ».}. Enfin
$\widetilde{\mathrm{C}}_{\mathrm{st}}(S)/\mathbb{D}_S \simeq
\mathrm{C}_{\mathrm{aff,st}}(n)$, l'ensemble des classes de systèmes simples
$(X, \Sigma, \Delta)$, et $\widetilde{\mathrm{C}}_{\mathrm{st}}(S)/\{1, a_S\}
\simeq \mathrm{C}^{!}_{\mathrm{aff,st}}(I_0)$.

\pagerange{9}{14}
\section*{Dessins, et une page étrangère (pages 9 à 14)}

\pagerange{9}{9}
\subsection*{Un arrangement de neuf courbes (page 9)}

La page est occupée par un grand arrangement tracé dans un disque. Ses
courbes portent, à leurs extrémités sur le bord, les étiquettes appariées
$\Delta_A, \Delta_{A'}, \Delta_B, \Delta_{B'}, \Delta_C, \Delta_{C'}$, et trois
autres, $D_A, D_B, D_C$, s'enroulent autour du centre ; les points
d'intersection s'appellent $A, B, C, A', B', C'$ au centre, puis $L, M, N$,
$L', M', N'$, $P, Q, R$, $P', Q', R'$. Chaque région porte son nombre de
côtés, et la page les compte : $7$ triangles, $24$ quadrilatères, $6$
pentagones, soit $37$ régions. Au-dessous, deux hexagones de sommets $A,
B', C, A', B, C'$.

Neuf pseudo-droites simples découpent le plan projectif en $1 + 9 \cdot 8 / 2
= 37$ régions, ce qui s'accorde avec le total. Mais la somme des nombres de
côtés, $3 \cdot 7 + 4 \cdot 24 + 5 \cdot 6 = 147$, est impaire, alors que pour
les faces d'une carte elle vaut deux fois le nombre d'arêtes ; les régions
« à l'infini », que la page ferme par un arc du bord, comptent sans doute ce
côté-là, ou un décompte est faux. On ne tranche pas sur la
transcription\footnote{La vérification de parité est de nous. Les mots
« triangles », « carrés », « pentagones » sont du transcripteur : sur la
page, chaque ligne du décompte est précédée d'un petit polygone dessiné.}.
Les mêmes lettres $A, A', B, C', L, N, P, Q, R, \ldots$ marquent la
construction de la page 13.

\pagerange{10}{10}
\subsection*{La factorisation de Stein (page 10)}

Cette page, sur une autre feuille, n'a aucun lien avec les pseudo-droites.
Pour une application $f : Z \to X$, notons $Z \to \widetilde{Z} \to X$ sa
\emph{factorisation de Stein} : en topologie, pour $f$ propre,
$\widetilde{Z}$ est l'espace des composantes connexes des fibres ; en
géométrie algébrique, $\widetilde{Z} = \mathrm{Spec}_X f_* \mathcal{O}_Z$.
Disons $f$ \emph{steinienne} si $\widetilde{Z} \simeq X$. La page énonce la
propriété universelle : si $Z' \to X'$ est steinienne et s'insère dans un
carré commutatif au-dessus de $Z \to X$, il existe une unique $X' \to
\widetilde{Z}$ telle que le carré

\begin{tikzcd}
Z \arrow[d] & Z' \arrow[l] \arrow[d] \\
\widetilde{Z} & X' \arrow[l]
\end{tikzcd}

commute. Autrement dit, dans la catégorie des flèches steiniennes $Z' \to X'$
au-dessus de $Z \to X$, l'objet $(Z \to \widetilde{Z})$ est \emph{final} :
de tout objet il part une unique flèche vers lui\footnote{La page dit « objet
initial ». Avec les flèches telles qu'elle les dessine --- de l'objet
$(Z' \to X')$ vers $(Z \to \widetilde{Z})$ ---, la propriété énoncée est
celle d'un objet final ; elle serait initiale dans la catégorie opposée.
Sous la flèche $X' \to \widetilde{Z}$, un mot souligné est illisible.}.

Il demande ensuite quand la factorisation commute au changement de base :
si le carré $Z' = Z \times_X X'$ est cartésien, les deux carrés de

\begin{tikzcd}
Z \arrow[d] & Z' \arrow[l] \arrow[d] \\
\widetilde{Z} \arrow[d] & \widetilde{Z}' \arrow[l] \arrow[d] \\
X & X' \arrow[l]
\end{tikzcd}

le sont-ils ? Il note deux conditions possibles, a) « $Z \to X$ compact ? »
et b) « $X' \to X$ localement $\mathcal{O}$-acyclique »\footnote{« compact »
et la lettre $\mathcal{O}$ sont des lectures incertaines.}, et s'arrête. Deux
réponses connues, que la page ne donne pas : en topologie, pour $f$ propre
entre espaces localement compacts, les fibres de $Z'$ étant celles de $Z$,
les deux carrés sont cartésiens pour tout changement de base --- c'est la
stabilité de la factorisation « monotone-légère » ; en géométrie algébrique,
pour $f$ quasi-compact et quasi-séparé, il suffit que $X' \to X$ soit plat
(changement de base plat pour $f_*\mathcal{O}_Z$)\footnote{Ces deux réponses
sont de nous, citées de mémoire ; la condition b) de la page va dans le
sens de la seconde.}.

\pagerange{11}{14}
\subsection*{Polyèdres et constructions (pages 11 à 14)}

La page 11 est une page de dessins au crayon : une sphère parcourue de trois
grands cercles, une pyramide et une bipyramide, un réseau de points, un
octaèdre inscrit dans un cube, un polyèdre à faces carrées et hexagonales
(un octaèdre tronqué ?) inscrit dans un cube, un arbre. Elle porte
\[
48 = 2 \cdot 4!, \qquad \frac{5!}{2 \cdot 4!} = \frac{5}{2}, \qquad
12 = \frac{8 \times 3 \times 2}{4}.
\]
Une sphère coupée par trois grands cercles en position générale \emph{est}
un octaèdre --- six sommets, douze arêtes, huit triangles --- et c'est le
revêtement double de trois droites du plan projectif : c'est l'objet des
pages 31 à 38, et $48 = 2 \cdot 4!$ est l'ordre de son groupe de symétries
$\mathfrak{S}_4 \times \{\pm 1\}$\footnote{Ce rapprochement est de nous ; la
page n'écrit rien autour de ces nombres.}.

La page 12 dessine un arrangement de six pseudo-droites, un pentagone central
entouré de cinq triangles, six régions marquées $4_1, \ldots, 4_6$. La page
14, cinq droites $D_0, \ldots, D_4$ formant une étoile, des demi-plans
hachurés, et une courbe --- une pseudo-droite --- qui coupe l'étoile.

La page 13 construit un arrangement de droites rectilignes : deux
horizontales, deux verticales et des obliques, portant les points $P, L, B,
C', N$ et $N', C, B', L', P'$, et des longueurs $a, b, a', b'$ sur des
segments. À droite :
\[
a + a' \leqslant 2, \qquad b + b' \geqslant 2, \qquad a' + b > 2,
\]
d'où $b > a$, puisque $a < 2 - a' < b$\footnote{Le signe des deux premières
inégalités a été repassé, et on ne sait pas s'il est large ou strict ; la
conclusion $b > a$ n'en dépend pas.}. Les lettres sont celles de la page 9,
et la page 9 fait s'enrouler trois de ses courbes : il est naturel de lire
la page 13 comme une tentative de réaliser l'arrangement de la page 9 par de
vraies droites, en suivant les contraintes métriques que cela impose. C'est
une conjecture de notre part ; elle rejoint la question de
l'\emph{étirabilité} posée page 60, et neuf est justement le plus petit
nombre de pseudo-droites d'un arrangement non étirable (Ringel, 1956, à
partir de la configuration de Pappus)\footnote{Rien sur les pages ne dit que
l'arrangement de la page 9 est celui-là, et on ne l'affirme pas.}.

\pagerange{16}{18}
\section*{II. Hyperplans réels et vecteurs de signes (pages 16 à 18)}

Soit $V$ un espace vectoriel réel de dimension finie et $H$ un hyperplan de
$V$. Les trois parties $H$, $V'_H$, $V''_H$ --- l'hyperplan et ses deux
demi-espaces ouverts --- forment une partition $D_H$ de $V$ en parties
convexes. C'est un ensemble à trois éléments dont un est distingué ; le
choix de l'un des deux demi-espaces comme positif l'identifie à
$\mathbf{F}_3 = \{-1, 0, 1\}$\footnote{La page dit « on peut regarder $D_H$
comme isomorphe à $\mathbf{F}_3$ » ; l'isomorphisme dépend du choix d'une
coorientation de $H$. Ce qui est canonique, c'est l'ordre défini plus bas,
où $0$ est le plus petit élément. La note marginale qui ouvre la page,
« Théorie analogue dans le cas affine, \ldots », est en majeure partie
illisible.}. Soit $\mathbf{S} = (V \smallsetminus \{0\})/\mathbf{R}^{*}_{+}$ la
sphère des directions\footnote{La page écrit $\mathbf{S} = V^{*}/\mathbf{R}^{*}_{+}$,
où $V^{*}$ ne peut être ici que $V \smallsetminus \{0\}$, et non le dual.}, et
$\varphi_H : \mathbf{S} \to D_H$ l'application qui envoie une direction sur la
partie qui la contient.

Pour une famille finie $(H_i)_{i \in I}$ d'hyperplans distincts, on obtient
\[
\varphi : \mathbf{S} \longrightarrow \mathbf{E} = \prod_{i \in I} D_{H_i},
\]
qui associe à chaque direction son \emph{vecteur de signes}. Les fibres non
vides $Z^{*}_\alpha = \varphi^{-1}(\alpha)$ forment une partition de
$\mathbf{S}$ en parties convexes --- leurs images inverses dans $V$ sont des
cônes convexes --- et $\Sigma = \varphi(\mathbf{S}) \subset \mathbf{E}$ est leur
ensemble d'indices. Il est stable par $-\mathrm{id}_{\mathbf{E}}$, et
$Z^{*}_{-\alpha} = -Z^{*}_\alpha$\footnote{Le mot « alvéoles ouvertes » est
celui de la page. Deux notes marginales, rattachées à « distincts » et à
« convexes », sont en grande partie illisibles ; on y lit deux fois « si
$\bigcap H_i = \{0\}$ ».}.

\textbf{Les adhérences.} Ordonnons chaque $D_{H_i}$ par $0 < \alpha'_i$ et $0 <
\alpha''_i$, les deux éléments non nuls étant incomparables : $\xi \leqslant
\eta$ si et seulement si $\xi = 0$ ou $\xi = \eta$. On munit $\mathbf{E}$ de
l'ordre produit et $\Sigma$ de l'ordre induit. Alors
\[
\overline{Z^{*}_\alpha} = \bigcup_{\beta \in \Sigma,\ \beta \leqslant \alpha}
Z^{*}_\beta :
\]
l'adhérence d'une strate est l'intersection des demi-espaces \emph{fermés}
et des hyperplans qui la définissent, et un point y est dans une strate
$\beta$ dont chaque signe est soit celui de $\alpha$, soit $0$\footnote{La
page écrit « On se propose de déterminer $\overline{Z}^{*}_\alpha$ pour
$i \in I$ », où il faut $\alpha \in \Sigma$. L'argument --- l'adhérence d'un
cône convexe ouvert dans son espace engendré, défini par des inégalités
strictes, est défini par les inégalités larges --- est développé par nous ;
la page l'énonce.}. Munissons $\Sigma$ de la topologie dont les fermés sont
les parties stables par minoration (la « topologie inférieure » de la page,
aujourd'hui topologie d'Alexandrov de l'ordre) ; alors $\beta \leqslant
\alpha$ équivaut à $\beta \in \overline{\{\alpha\}}$, et pour toute partie $A
\subset \Sigma$, en posant $Z^{*}_A = \bigcup_{\alpha \in A} Z^{*}_\alpha$,
\[
\overline{Z^{*}_A} = Z^{*}_{\overline{A}},
\]
puisque $\Sigma$ est fini et que l'adhérence d'une réunion finie est la
réunion des adhérences. L'ensemble ordonné $\Sigma$ est donc l'ensemble
d'indices d'une stratification de la sphère $\mathbf{S}$ en strates convexes,
qui satisfait l'axiome de frontière : l'adhérence d'une strate est réunion de
strates. Lorsque $\bigcap_i H_i = \{0\}$, les strates sont des cellules
ouvertes et l'on obtient une décomposition cellulaire régulière de la
sphère ; sinon, la strate de signe nul est la sphère de $\bigcap_i H_i$, et
l'on se ramène à ce cas en passant au quotient par
$\bigcap_i H_i$\footnote{Cette dernière remarque est de nous ; c'est sans
doute ce que vise l'hypothèse « si $\bigcap H_i = \{0\}$ » des notes
marginales.}.

Dans le vocabulaire actuel, $\Sigma$ est l'ensemble des \emph{covecteurs}
non nuls du \emph{matroïde orienté} réalisé par les $H_i$, et
$(\Sigma, \leqslant)$ est le \emph{poset des faces} de l'arrangement. Le pont
avec les pseudo-droites, que ce dossier cherche partout ailleurs, est le
\emph{théorème de représentation topologique} de Folkman et Lawrence (1978) :
les matroïdes orientés de rang $3$ sont exactement les arrangements de
pseudo-droites (de pseudo-cercles sur la sphère)\footnote{Les noms sont de
nous ; les pages n'en citent aucun.}.

\textbf{Les bornes inférieures.} Un paragraphe barré dit que $\Sigma$ n'est
pas en général un ensemble où deux éléments minorés ont une borne inférieure,
avec l'exemple de trois plans de $\mathbf{R}^3$ passant par une même droite :
deux fuseaux opposés de la sphère ont pour minorants communs les deux points
$\pm p$ de cette droite, et aucune borne inférieure. Le paragraphe qui suit,
non barré, affirme que $\Sigma$ est une partie ouverte de $\mathbf{E}$,
c'est-à-dire stable par majoration, et en déduit que toute partie minorée de
$\Sigma$ y a une borne inférieure.

La première affirmation est fausse en général, et le contre-exemple est
celui que la page vient de barrer. Ce qui est vrai : \emph{$\Sigma$ est
stable par majoration dans $\mathbf{E}$ si et seulement si, pour tout $x
\neq 0$, les formes linéaires des hyperplans qui contiennent $x$ sont
linéairement indépendantes} --- par exemple si les $H_i$ sont en position
générale. Si elles le sont, on peut donner aux $\ell_i(x + \varepsilon v)$
les signes qu'on veut ; si une relation $\sum c_i \ell_i = 0$ les lie, le
vecteur de signes $(\operatorname{sgn} c_i)$ est interdit. Sous cette
hypothèse, la conclusion suit : dans $\mathbf{E}$, toute partie non vide a
une borne inférieure (coordonnée par coordonnée), et si la partie est minorée
dans $\Sigma$, cette borne majore un élément de $\Sigma$, donc est dans
$\Sigma$\footnote{L'énoncé corrigé et son argument sont de nous. La fin de
la page est d'une lecture très incertaine ; « Inf » y est écrit « I--f ».}.

\pagerange{20}{28}
\section*{III. Les arrangements cruciaux (pages 20 à 28)}

\pagerange{20}{20}
\subsection*{Une première liste (page 20)}

La page 20 dessine quatre arrangements : (1) quatre droites en position
générale, de groupe $\mathfrak{S}_4$ ; (2) cinq droites en position
générale, de groupe diédral $D_5$ ; (3) cinq droites « en position
particulière » ; (4) six droites, les côtés d'un triangle et ses trois
médianes concourantes. Elle énonce : un arrangement de pseudo-droites sans
pseudo-droite bigonale, s'il ne se réduit pas à quatre pseudo-droites (cas
(1)), contient une configuration de l'un des types (2), (3), (4) ; et ces
types sont « non spéciaux » : réalisés en arrangements de dimension $2$,
l'intersection de deux faces fermées y est vide, un sommet ou une arête
fermée\footnote{Le début de l'énoncé est d'une lecture difficile
(« Tout arr.\ \ldots{} $1$-dim.\ de ps.\ droites \ldots{} pouvant être de
ps.\ droites bigones ») ; on lui donne l'hypothèse « sans bigone » de la page
24, qui reprend le même énoncé. La parenthèse ouverte avant « de
ps.\ droites » n'est pas refermée.}. Un NB ajoute : s'il ne contient ni (2) ni
(3), il est du type (4) ou d'un type (5), sept pseudo-droites formées d'un
quadrilatère et de ses trois diagonales, « qui mérite d'être étudié ». La
page s'arrête sur « Si elle ne contient », biffé. Les pages 22 et 24 donnent
des noms à tous ces types ; on les retrouve ci-dessous, avec la correction
que le NB demande.

\pagerange{22}{22}
\subsection*{La liste des types (page 22)}

On cherche les « cas cruciaux » d'arrangements de $N \geqslant 4$
pseudo-droites, le cas $N = 3$ étant connu. Le nombre de sommets de chacun,
que la page note en marge, est rappelé entre parenthèses.

\emph{1\textsuperscript{o}) Il existe une pseudo-droite bigonale.}
\begin{itemize}
  \item $\mathrm{I}_N$ : un faisceau de $N - 1 \geqslant 3$ pseudo-droites
  concourantes et une sécante ($N$ sommets)\footnote{La marge porte
  « $N+1$ sommets ». Le centre du faisceau et les $N-1$ points de la sécante
  font $N$ sommets.}. Les pseudo-droites du faisceau sont bigonales.
  \item $\mathrm{I}_{p,q}$ : deux faisceaux de $p$ et $q$ pseudo-droites,
  $p \geqslant q \geqslant 2$, et la pseudo-droite qui joint leurs deux
  centres, $p + q + 1 = N \geqslant 5$ ($pq + 2$ sommets). Celle-ci est
  l'unique pseudo-droite bigonale.
\end{itemize}
Ces deux types épuisent le cas : si $D$ est bigonale, de sommets $x$ et $y$,
toute autre pseudo-droite coupe $D$ en $x$ ou en $y$, et l'on a deux
faisceaux de centres $x$ et $y$ ; s'il n'y a qu'une pseudo-droite par $y$
hors de $D$, c'est le type $\mathrm{I}_N$\footnote{Cette vérification est de
nous ; la page donne la liste sans dire qu'elle est complète.}.

\emph{2\textsuperscript{o}) Il n'existe pas de pseudo-droite bigonale.}
\begin{itemize}
  \item $\mathrm{II}'_{p,q}$ : deux faisceaux de $p$ et $q$ pseudo-droites,
  $p \geqslant q \geqslant 2$, $p + q = N$, sans la pseudo-droite qui joindrait
  leurs centres ($pq + 2$ sommets).
  \item $\mathrm{II}_4$ : quatre pseudo-droites en position générale ($6$
  sommets).
  \item $\mathrm{II}_5$ : cinq pseudo-droites en position générale ($10$
  sommets).
\end{itemize}
Le cas « minimal » $p = q = 2$ de $\mathrm{II}'_{p,q}$ est justement
$\mathrm{II}_4$ : deux paires de pseudo-droites se coupant en deux points
distincts, sans troisième point de concours\footnote{L'identification
$\mathrm{II}'_{2,2} = \mathrm{II}_4$ est de nous ; la page écrit que ce cas
minimal est « justiciable \ldots », la fin étant illisible.}.

\pagerange{24}{24}
\subsection*{Le théorème (page 24)}

Deux autres types entrent ici.
\begin{itemize}
  \item $\mathrm{III}_6$, l'arrangement « tétraédral » : les six
  pseudo-droites qui joignent deux à deux quatre points, dont trois jamais
  alignés --- le \emph{quadrangle complet} ($4 + 3 = 7$ sommets : quatre
  points triples et trois points doubles, les points diagonaux). C'est la
  configuration (4) de la page 20 : un triangle et ses trois médianes.
  \item $\mathrm{IV}_7$ : un \emph{quadrilatère complet} --- quatre
  pseudo-droites en position générale et leurs six sommets --- et ses trois
  diagonales, qui joignent les sommets opposés ($6 + 3 = 9$ sommets : six
  points triples et les trois points diagonaux). C'est la configuration (5)
  de la page 20.
\end{itemize}

\emph{Théorème} (énoncé, non démontré). Soit $K$ un arrangement de
pseudo-droites tel qu'il n'existe pas deux sommets par l'un desquels passe
chaque pseudo-droite --- c'est-à-dire qui n'est d'aucun des types
$\mathrm{I}_N$, $\mathrm{I}_{p,q}$, $\mathrm{II}'_{p,q}$ --- et qui ne contient
pas de sous-arrangement de type $\mathrm{II}_5$. Alors $K$ est de type
$\mathrm{III}_6$ ou $\mathrm{IV}_7$.

La page en tire une forme encadrée : \emph{tout arrangement sans bigone est
de type $\mathrm{II}_4$ ou contient un sous-arrangement de l'un des types}
\[
\mathrm{II}'_{3,2}, \qquad \mathrm{II}_5, \qquad \mathrm{III}_6,
\]
et ces trois types sont « latticiels » : les facettes fermées de
l'arrangement de dimension $2$ correspondant forment un treillis,
l'intersection de deux facettes fermées non vides étant une
facette\footnote{Le premier type est écrit $\mathrm{I}'_{3,2}$ ici et
$\mathrm{II}_{3,2}$ aux pages 26 et 28, et n'est défini nulle part. On le lit
$\mathrm{II}'_{3,2}$ --- trois pseudo-droites par un point, deux par un
autre, sans droite de jonction ---, ce qui est notre lecture : c'est le seul
arrangement de cinq pseudo-droites sans bigone autre que $\mathrm{II}_5$, et
c'est ce que demande la page 28, où un sous-arrangement de cinq
pseudo-droites sans bigone est dit « du type $\mathrm{II}_5$ ou du type
$\mathrm{II}_{3,2}$ ». C'est aussi la configuration (3) de la page 20.}. La
forme encadrée découle du théorème : si $K$ est couvert par deux faisceaux
sans bigone, c'est un $\mathrm{II}'_{p,q}$, qui contient
$\mathrm{II}'_{3,2}$ dès que $p \geqslant 3$ et vaut $\mathrm{II}_4$ sinon ; et
$\mathrm{IV}_7$ contient $\mathrm{III}_6$ (deux diagonales et les quatre côtés
forment le quadrangle complet de quatre sommets du quadrilatère).

La page ajoute entre parenthèses que s'il ne contient ni $\mathrm{II}'_{3,2}$
ni $\mathrm{II}_5$, il est de type $\mathrm{III}_6$ ou $\mathrm{IV}_7$ ; le NB
de la page 20 dit la même chose. Cela ne peut pas être juste tel quel :
$\mathrm{IV}_7$ contient un $\mathrm{II}'_{3,2}$. Si $A, B, C, D$ sont les
sommets d'un quadrangle, $E = AB \cap CD$, $F = AD \cap BC$, les sept
pseudo-droites de $\mathrm{IV}_7$ sont $AB$, $BC$, $CD$, $DA$, $AC$, $BD$,
$EF$, et les cinq pseudo-droites $AB$, $AD$, $AC$, $BD$, $EF$ n'ont qu'un
point triple, $A$. Le théorème admis, la conclusion correcte est : \emph{s'il
ne contient ni $\mathrm{II}'_{3,2}$ ni $\mathrm{II}_5$, $K$ est de type
$\mathrm{II}_4$ ou $\mathrm{III}_6$}\footnote{La correction et le
contre-exemple sont de nous. On vérifie aussi que $\mathrm{III}_6$ ne
contient ni $\mathrm{II}_5$ ni $\mathrm{II}'_{3,2}$ : ôter une pseudo-droite
du quadrangle complet laisse deux points triples. La forme encadrée n'est pas
touchée.}.

\pagerange{26}{26}
\subsection*{Le cas latticiel (page 26)}

On est dans le \emph{cas latticiel} si l'intersection de deux facettes
fermées est vide ou une facette fermée. Grothendieck énonce : \emph{on est
dans le cas latticiel si et seulement s'il n'y a pas de pseudo-droite
bigonale et que l'arrangement n'est pas de type $\mathrm{II}_4$}\footnote{La
page écrit « b) pas de type $\mathrm{II}_4$ ». On lit « $K$ n'est pas de type
$\mathrm{II}_4$ » plutôt que « ne contient pas » : $\mathrm{II}_5$ contient
$\mathrm{II}_4$ et la page 24 le dit latticiel. Dans $\mathrm{II}_4$, deux
des trois quadrilatères ont en commun deux sommets opposés, qui ne forment
pas une facette.}.

\emph{Proposition} (énoncée). Soit $\mathcal{X}$ un arrangement de dimension
$2$ de $N$ pseudo-droites $(D_i)_{i \in I}$, et $D$ une $(N+1)$-ième
pseudo-droite, qui ne soit pas bigonale dans l'arrangement $\mathcal{X}'$
obtenu en l'ajoutant. Soit $S$ l'ensemble des sommets de $\mathcal{X}'$ sur
$D$, c'est-à-dire des points $D \cap D_i$. Deux points $s, s' \in S$ sont
adjacents sur $D$ si et seulement s'il existe une facette fermée $f$ de
$\mathcal{X}$ qui les contient tous deux ; elle est alors unique, et $D \cap
f$ est l'arête de $\mathcal{X}'$ qui joint $s$ et $s'$\footnote{La fin de
l'énoncé, autour de « $D \cap f$ » et d'une courbe notée $T$ sur le dessin,
est d'une lecture incertaine ; on retient ce que dit la dernière ligne :
« c'est aussi l'arête sur $D$ joignant $s$ et $s'$ ». Un premier énoncé du
lemme de la page 28, barré, suit sur la page.}.

C'est l'outil de la récurrence : si l'on ajoute les pseudo-droites une à une
sans jamais créer de bigone, la position de chaque nouvelle pseudo-droite
dans les faces de l'arrangement précédent se lit sur le squelette.

\pagerange{28}{28}
\subsection*{Le lemme d'ordonnancement (page 28)}

\emph{Lemme} (énoncé ; seule une « Idée » de démonstration est écrite).
Soit $K$ un arrangement de dimension $1$ de $N \geqslant 5$ pseudo-droites,
sans pseudo-droite bigonale. On peut numéroter ses pseudo-droites $D_1,
\ldots, D_N$ de sorte que, $K_i$ désignant l'arrangement formé de $D_1,
\ldots, D_i$,
\begin{enumerate}
  \item[1\textsuperscript{o})] $D_1, \ldots, D_4$ forment un arrangement de type
  $\mathrm{II}_4$ (ou bien $D_1, \ldots, D_6$ un arrangement de type
  $\mathrm{III}_6$) ;
  \item[2\textsuperscript{o})] pour $i \geqslant 5$ (respectivement $i
  \geqslant 7$), $D_i$ n'est pas bigonale dans $K_i$.
\end{enumerate}
L'idée procède par cas, que la page dit mutuellement exclusifs : si l'on
trouve un $K_5$ de type $\mathrm{II}_5$, on continue d'ajouter des
pseudo-droites ; si l'on trouve un $K_5$ de type $\mathrm{II}'_{3,2}$, une
pseudo-droite ajoutée n'est jamais bigonale ; si $K$ contient un
$\mathrm{III}_6$, on part de lui. Le point de départ est donc toujours
$\mathrm{II}_4$ ou $\mathrm{III}_6$\footnote{Ce passage est écrit très vite,
et seules les formules s'en lisent sûrement ; un « $\mathrm{III}_7$ » y
paraît deux fois sans être défini. On ne reconstitue pas l'argument.}.

\pagerange{31}{38}
\section*{IV. Trois pseudo-droites : cubes, octaèdres et leurs formes gauches (pages 31 à 38)}

Ces pages, d'une écriture posée et numérotées, traitent à fond le cas $N = 3$
que la page 22 disait « bien étudié ». Trois pseudo-droites en position
générale découpent le plan projectif en $3$ sommets, $6$ arêtes et $4$
triangles : c'est la moitié d'un octaèdre, l'\emph{hémi-octaèdre}, et son
revêtement double sur la sphère est l'octaèdre lui-même.

\pagerange{31}{32}
\subsection*{Cube et octaèdre combinatoires (pages 31 et 32)}

Soit $I$ un ensemble à trois éléments et $(\Phi_i)_{i \in I}$ une famille
d'ensembles à deux éléments ; posons $\Phi = \coprod_i \Phi_i$, à six
éléments, au-dessus de $I$. Le \emph{cube combinatoire} de dimension $3$
associé a pour faces les éléments de $\Phi$, pour arêtes les sections
partielles de $\Phi$ au-dessus d'une partie à deux éléments de $I$, pour
sommets les sections totales :
\[
C_2 \simeq \Phi, \qquad C_1 \simeq \{A \in \mathfrak{P}_2(\Phi) \mid A \to I
\text{ injectif}\}, \qquad C_0 \simeq \{A \in \mathfrak{P}_3(\Phi) \mid A \to I
\text{ bijectif}\},
\]
l'incidence étant l'inclusion renversée. Les éléments de $I$ correspondent
aux paires de faces opposées. L'\emph{octaèdre combinatoire}, dual du cube,
a les mêmes ensembles en dimensions renversées, $O_0 = C_2$, $O_1 = C_1$,
$O_2 = C_0$, l'incidence étant l'inclusion.

L'automorphisme involutif $\underline{a}$ de $\Phi$, identique sur $I$ et
échangeant les deux éléments de chaque $\Phi_i$, induit l'\emph{antipodie}
du cube et de l'octaèdre ; il est central. Les quotients
\[
\widetilde{C}_k = C_k/\{1, \underline{a}\}, \qquad
\widetilde{O}_k = O_k/\{1, \underline{a}\}
\]
sont le cube et l'octaèdre \emph{gauches}, ou projectifs --- aujourd'hui
hémi-cube et hémi-octaèdre. L'octaèdre gauche a $3$ sommets, $6$ arêtes et
$4$ faces : c'est la décomposition du plan projectif réel par trois droites
indexées par $I$. À $i \in I$ correspond la « droite » $D_i$ formée des deux
arêtes projectives images des quatre arêtes $A \in \mathfrak{P}_2(\Phi)$
au-dessus de $I \smallsetminus \{i\}$\footnote{Sur la page, « gauche (ou
projectif) » est ajouté en interligne, et plus bas « gauche » remplace
« projectif » biffé ; le mot se lit sûrement page 37. La fin de la page 32 est
d'une lecture incertaine.}.

\pagerange{33}{35}
\subsection*{Orientations (pages 33 à 35)}

Les foncteurs $C \mapsto \widetilde{C}$ et $O \mapsto \widetilde{O}$ ne sont
pas des équivalences : l'antipodie est un automorphisme non trivial de $C$
et de $O$ qui induit l'identité de $\widetilde{C}$ et $\widetilde{O}$. Ils
sont \emph{pleins} --- tout automorphisme du quotient se relève ---, mais
\emph{pas fidèles} : sur les groupes d'automorphismes, c'est la surjection
$\mathfrak{S}_4 \times \{\pm 1\} \to \mathfrak{S}_4$ de noyau $\{1,
\underline{a}\}$\footnote{La page écrit qu'ils « ne sont pas pleinement
fidèles mais \ldots{} fidèles, car l'antipodisme $\underline{a}$ est un
automorphisme non trivial de $C$, $O$ et induit l'automorphisme trivial de
$\widetilde{C}$, $\widetilde{O}$ » ; l'argument donné prouve le contraire de
la fidélité, et la page 38 dit bien que $\mathrm{oct} \to
\mathrm{oct.\ gauches}$ est « surjectif sur les isom. ». La lecture « sont » est
incertaine.}. On obtient en revanche des équivalences de catégories en
prenant des cubes et des octaèdres \emph{orientés} :

(5)

\begin{tikzcd}
\text{cubes comb.\ orientés} \arrow[r, "\approx"] \arrow[d, "\approx"'] &
\text{cubes comb.\ gauches} \arrow[d, "\approx"] \\
\text{octaèdres comb.\ orientés} \arrow[r, "\approx"] &
\text{octaèdres comb.\ gauches}
\end{tikzcd}

car le groupe des rotations du cube, $\mathfrak{S}_4$, ne contient pas
l'antipodie (qui renverse l'orientation en dimension $3$) et s'envoie
isomorphiquement sur le groupe de l'hémi-cube. Inversement, à partir de
$\widetilde{C}$ ou $\widetilde{O}$, on retrouve $C$ ou $O$ orienté comme
revêtement d'orientation.

L'ensemble des orientations du cube, ou de l'octaèdre, défini par $(\Phi_i)$
est
\[
(6) \qquad \mathrm{or}(C) \simeq \mathrm{or}(O) \simeq \Bigl(\bigwedge_{i \in I}
\Phi_i\Bigr) \wedge \mathrm{or}(I) :
\]
une orientation du cube, réalisé comme le cube $[-1,1]^{I}$, est
donnée par un choix d'un élément de chaque $\Phi_i$ (un sens sur chaque axe)
et d'un ordre circulaire sur $I$ (l'ordre des axes), et changer l'un des
choix renverse l'orientation. Une orientation est donc un isomorphisme
\[
(6\ \text{bis}) \qquad \Bigl(\bigwedge_{i \in I} \Phi_i\Bigr) \wedge \mathrm{or}(I)
\simeq \underline{1}, \quad \text{ou encore} \quad \bigwedge_{i \in I} \Phi_i
\simeq \mathrm{or}(I).
\]

Dans l'octaèdre gauche, la droite $D_i$ porte deux sommets et deux arêtes.
Les sommets de $\widetilde{O}$ sont les paires $\Phi_j$, donc indexés par $I$,
et $D_i$ passe par les deux sommets autres que $i$ :
\[
(8) \qquad S_i \simeq I \smallsetminus \{i\}.
\]
Ses deux arêtes sont les deux classes, modulo l'antipodie, d'arêtes $\{x_j,
x_k\}$ du cube au-dessus de $\{j, k\} = I \smallsetminus \{i\}$ ; une telle
classe est une bijection $\Phi_j \simeq \Phi_k$, d'où
\[
(7) \qquad A_i = \bigwedge_{\alpha \in I \smallsetminus \{i\}} \Phi_\alpha .
\]
Chaque $D_i$ est donc un \emph{bigone} du graphe de $\widetilde{O}$. Orienter
un bigone, c'est dire de quel sommet part chacune de ses deux arêtes ;
changer le sommet ou l'arête renverse l'orientation, si bien que
l'ensemble des deux orientations de $D_i$ est
\[
(9) \qquad \Omega_i \simeq S_i \wedge A_i .
\]
Chaque $\Phi_\alpha$ figurant dans deux des $A_i$,
\[
(10) \qquad \bigwedge_{i \in I} A_i \simeq \bigwedge_{\alpha \in I}
\Phi_\alpha^{\wedge 2} \simeq \underline{1},
\]
et donc
\[
(11) \qquad \bigwedge_{i \in I} \Omega_i \simeq \bigwedge_{i \in I} S_i \simeq
\bigwedge_{i \in I} (I \smallsetminus \{i\}) \simeq \mathrm{or}(I),
\]
le dernier isomorphisme étant l'isomorphisme canonique évident pour un
ensemble à trois éléments : à un ordre circulaire $\sigma$ on associe le
choix, dans chaque $I \smallsetminus \{i\}$, de l'élément $\sigma(i)$ ; renverser
$\sigma$ change les trois choix\footnote{L'explicitation de cet isomorphisme
est de nous ; la page dit « isom.\ can.\ évident associé à l'ens.\ à $3$
éléments $I$ ».}. Pour trois bigones quelconques $(S_i, A_i)$ avec $S_i = I
\smallsetminus \{i\}$, se donner l'isomorphisme (10) revient donc à se donner
\[
(11\ \text{bis}) \qquad \bigwedge_{i \in I} \Omega_i \simeq \mathrm{or}(I), \qquad
\text{soit} \qquad (12) \qquad \Bigl(\bigwedge_{i \in I} \Omega_i\Bigr) \wedge
\mathrm{or}(I) \simeq \underline{1},
\]
c'est-à-dire une orientation du cube défini par la famille $(\Omega_i)$.

Une fois le cube orienté par (6 bis), les $\Omega_i$ et les $\Phi_i$ se
déterminent mutuellement :
\[
(13) \qquad \Omega_i \simeq \Phi_i \wedge \bigwedge_{\alpha \in I
\smallsetminus \{i\}} (I \smallsetminus \{\alpha\}), \qquad \Phi_i \simeq \Omega_i
\wedge \bigwedge_{\alpha \in I \smallsetminus \{i\}} (I \smallsetminus
\{\alpha\}),
\]
d'où $\bigwedge \Omega_i \simeq \bigwedge \Phi_i$, et (14) : les orientations
des octaèdres définis par $(\Omega_i)$ et par $(\Phi_i)$ se correspondent.

\pagerange{35}{37}
\subsection*{Le graphe octaédral gauche (pages 35 à 37)}

\emph{Proposition} (« En résumé »). Considérons la carte de l'octaèdre gauche
--- trois sommets, six arêtes, quatre faces triangulaires ---, c'est-à-dire
un arrangement combinatoire de trois pseudo-droites $D_i$, $i \in I$, dans un
plan projectif. Soit $K$ son squelette de dimension $1$, muni de ses trois
sous-complexes $D_i$, qui sont des bigones, de sommets $S_i \simeq I
\smallsetminus \{i\}$ et d'arêtes $A_i$ ; on l'appelle un \emph{graphe
octaédral gauche}. (Les sommets de $K$ sont en bijection avec $I$, en
associant à $i$ l'unique sommet qui n'est pas sur $D_i$.) Posons
\[
\Omega_i = S_i \wedge A_i, \qquad \Phi_i = \Omega_i \wedge \bigwedge_{\alpha
\in I \smallsetminus \{i\}} (I \smallsetminus \{\alpha\}) \simeq A_i \wedge
\bigwedge_{\alpha \in I} (I \smallsetminus \{\alpha\}).
\]
On a des isomorphismes canoniques, définis directement à partir de $K$ et
des trois $D_i$,
\[
(15) \qquad \bigwedge_{i \in I} \Omega_i \simeq \bigwedge_{i \in I} \Phi_i
\simeq \Bigl(\bigwedge_{i \in I} A_i\Bigr) \wedge \mathrm{or}(I),
\]
qu'on peut réécrire
\[
(15\ \text{bis}) \qquad \Bigl(\bigwedge_{i} \Omega_i\Bigr) \wedge \mathrm{or}(I)
\simeq \Bigl(\bigwedge_{i} \Phi_i\Bigr) \wedge \mathrm{or}(I) \simeq
\bigwedge_{i} A_i ,
\]
les deux premiers termes étant les orientations des octaèdres définis par
$(\Omega_i)$ et par $(\Phi_i)$. Du fait que $K$ provient de la carte, le
torseur (15 bis) reçoit une \emph{trivialisation canonique}, celle de (10).
Appelons \emph{orientation} d'un graphe octaédral gauche une trivialisation
$\bigwedge_{i} A_i \simeq \underline{1}$. Précisément, le foncteur
\[
(16) \qquad (\text{octaèdres gauches}) \longrightarrow (\text{graphes
octaédraux gauches orientés})
\]
est une \emph{équivalence} de catégories. On le vérifie sur les groupes : le
graphe octaédral gauche seul a pour groupe d'automorphismes
$\mathfrak{S}_3 \ltimes (\mathbf{Z}/2\mathbf{Z})^3$, d'ordre $48$ (permuter les
sommets, et échanger indépendamment les deux arêtes de chaque bigone) ; la
trivialisation ne laisse que les échanges en nombre pair, $\mathfrak{S}_3
\ltimes (\mathbf{Z}/2\mathbf{Z})^2 \simeq \mathfrak{S}_4$, le groupe de
l'hémi-octaèdre\footnote{Cette vérification est de nous. La
trivialisation a aussi une lecture directe : parmi les huit triples
d'arêtes, une sur chaque $D_i$, elle distingue une classe de quatre, et ce
sont les bords des quatre faces triangulaires. C'est en ce sens que les
pages 46, 52 et 54 parlent de triples d'arêtes « compatibles avec
l'orientation ». La note marginale de la page 37, qui glose l'équivalence
comme l'ajout d'« une orientation des gr.\ octaédraux gauches », est d'une
lecture incertaine.}. 

\pagerange{37}{38}
\subsection*{Le cube de foncteurs (pages 37 et 38)}

Grothendieck résume les relations entre octaèdres et graphes octaédraux,
gauches ou ordinaires, orientés ou non, par un diagramme commutatif en forme
de cube :

(17)

\begin{tikzcd}[column sep=small, row sep=small]
\text{oct}^{+} \arrow[rr, "\approx"] \arrow[dr, "\approx"] \arrow[dd] & &
\text{oct.\ g.} \arrow[dr, "\approx"] \arrow[dd, no head, "=" description] & \\
& \text{gr.\ oct}^{+} \arrow[rr, "\approx"] \arrow[dd] & &
\text{gr.\ oct.\ g.}^{+} \arrow[dd] \\
\text{oct} \arrow[rr] \arrow[dr, "\approx"] & & \text{oct.\ g.} \arrow[dr] & \\
& \text{gr.\ oct} \arrow[rr] & & \text{gr.\ oct.\ g.}
\end{tikzcd}

où « g. » abrège « gauche » et l'exposant $+$ désigne les objets orientés.
La face supérieure (objets orientés) est faite d'équivalences. Les quatre
flèches verticales, qui oublient l'orientation, sont fidèles ; celle de
droite à l'arrière est l'identité. La face arrière est formée des
configurations de dimension $2$, la face avant des configurations de
dimension $1$ associées ; les quatre flèches de l'arrière vers l'avant sont
des équivalences, sauf $\text{oct.\ g.} \to \text{gr.\ oct.\ g.}$, qui est
fidèle sans être plein, et dont l'énoncé (16) dit comment le rendre
plein\footnote{La page marque certaines flèches d'un petit crochet, qu'on ne
sait pas interpréter, et la flèche horizontale du milieu à l'arrière ne porte
pas de $\approx$. On n'a pas reproduit les crochets.}.

Sur la face inférieure, $\text{oct} \to \text{oct.\ g.}$ est pleine et non
fidèle, comme on l'a vu. La page affirme que $\text{gr.\ oct} \to
\text{gr.\ oct.\ g.}$ est une équivalence ; ce n'est pas possible si le cube commute.
En effet la composée $\text{oct} \to \text{oct.\ g.} \to \text{gr.\ oct.\ g.}$
induit sur les groupes $\mathfrak{S}_4 \times \{\pm 1\} \to \mathfrak{S}_4
\hookrightarrow \mathfrak{S}_3 \ltimes (\mathbf{Z}/2\mathbf{Z})^3$, de noyau
$\{1, \underline{a}\}$ et d'image d'indice $2$ ; puisque $\text{oct} \to
\text{gr.\ oct}$ est une équivalence, la flèche $\text{gr.\ oct} \to
\text{gr.\ oct.\ g.}$ n'est ni fidèle ni pleine. C'est d'ailleurs ce qu'on voit
directement : le passage au quotient par l'antipodie la tue\footnote{La
page dit : « les deux flèches $\text{oct} \to \text{gr.\ oct.} \to
\text{gr.\ oct g.}$ sont des équivalences » et, plus bas, que les quatre
flèches de passage du sphérique au gauche sont des équivalences « à
l'exception de $\text{oct} \to \text{oct.\ gauches}$ ». La correction est de
nous. Les groupes abstraits de $\text{gr.\ oct}$ et de $\text{gr.\ oct.\ g.}$
sont isomorphes, d'ordre $48$ tous deux, ce qui explique peut-être
l'affirmation ; mais le foncteur naturel n'induit pas cet isomorphisme.}. La
marge de la page 38 relie les groupes $\mathfrak{S}_4 \times \{\pm 1\}$ et
$\mathfrak{S}_4$ aux deux premières lignes de l'explication.

\pagerange{39}{40}
\section*{V. Arrangements de dimension 1 orientés (pages 39 et 40)}

\pagerange{39}{39}
\subsection*{Le théorème de pleine fidélité (page 39)}

Un arrangement de dimension $1$ $(K, (D_i)_{i \in I})$ est défini à la page
40 (ci-dessous). Une \emph{orientation} de cet arrangement est la donnée, pour tout $J
\in \mathfrak{P}_3(I)$ dont les trois pseudo-droites ne sont pas concourantes,
d'une orientation $\omega_J$ du graphe octaédral gauche $K_J$, au sens de
(16).

Soit $\mathcal{X}$ une carte qui est un arrangement de pseudo-droites. Son
squelette $K$ est muni de la famille $(D_i)$ de ses pseudo-droites, qui en fait
un arrangement de dimension $1$, et, par (16) appliqué à chaque triple, d'une
orientation $\omega = (\omega_J)_J$. D'où un foncteur
\[
(18) \qquad \mathcal{X} \longmapsto (K, (D_i)_{i \in I}, \omega)
\]
de la catégorie des arrangements combinatoires de pseudo-droites vers celle
des arrangements de dimension $1$ orientés.

\emph{Théorème} (énoncé page 39). \emph{Le foncteur (18) est pleinement
fidèle} : pour deux arrangements $\mathcal{X}$, $\mathcal{X}'$,
\[
\mathrm{Isom}(\mathcal{X}, \mathcal{X}') \xrightarrow{\ \sim\ }
\mathrm{Isom}\bigl((K, (D_i), \omega), (K', (D'_{i'}), \omega')\bigr).
\]
La démonstration occupe les pages 41 à 49 ; Grothendieck dit à la fin ne pas
l'avoir vérifiée dans un cas (voir plus bas)\footnote{Dans la marge de la
page 39, trois lignes barrées, où l'on devine « bigone », sont trop
incertaines pour être données.}.

\pagerange{40}{40}
\subsection*{Les axiomes (page 40)}

La page 40 donne après coup la définition que la page 39 suppose. Un
\emph{arrangement de dimension $1$ de pseudo-droites} est la donnée d'un
$1$-complexe combinatoire fini $K$ et d'une famille finie $(D_i)_{i \in I}$ de
sous-complexes polygonaux, tels que
\begin{itemize}
  \item toute arête de $K$ appartient à un $D_i$ et un seul ;
  \item pour $i \neq j$, $D_i$ et $D_j$ se rencontrent en un sommet et un seul ;
  \item tout sommet de $K$ appartient à au moins deux des $D_i$.
\end{itemize}
Supposons que l'intersection de tous les $D_i$ soit vide. Pour $J \subset I$,
$K_J$ est la réunion des $D_i$, $i \in J$, où l'on ne garde comme sommets que
les intersections de ces $D_i$. Si $J$ a trois éléments et $\bigcap_{i \in J}
D_i = \emptyset$, $K_J$ est un graphe octaédral gauche --- trois bigones
deux à deux sécants en un seul sommet ---, et l'on peut parler de ses
orientations\footnote{La page écrit deux fois « graphe octogonal gauche » ;
le contexte et la page 37 demandent « octaédral ». Elle est très reprise, et
une phrase commencée, « Il existe alors une partie $I \subset J$ de cardinal
\ldots », reste en suspens. La note marginale n'est lisible que par
fragments.}.

\pagerange{41}{49}
\section*{VI. Reconstruire la carte (pages 41 à 49)}

\pagerange{41}{42}
\subsection*{Ce qu'il faut reconstruire (pages 41 et 42)}

Pour démontrer le théorème, on montre comment reconstruire la carte
$\mathcal{X}$ à partir de $(K, (D_i), \omega)$. Une carte sur une surface
fermée, orientable ou non, dont les sommets sont d'ordre $\geqslant 3$, est
déterminée par son graphe et par les données suivantes (le « théorème
général des cartes » de la page) :
\begin{enumerate}
  \item[1\textsuperscript{o})] pour chaque sommet $s \in S$ (l'ensemble des sommets
  de $K$), la paire
  \[
  (19) \qquad \vec{\mathfrak{Z}}_s \subset \mathcal{C}_{\vec{A}_s}
  \]
  des deux ordres circulaires, inverses l'un de l'autre, dans lesquels les
  arêtes d'origine $s$ se succèdent autour de $s$ ; l'ensemble $\vec{A}_s$
  de ces arêtes s'identifie à l'ensemble $\vec{I}_s$ des pseudo-droites
  orientées $\vec{D}_i$ passant par $s$ ;
  \item[2\textsuperscript{o})] pour chaque arête orientée $\vec{a}$, d'origine $s$ et
  d'extrémité $s'$, une bijection
  \[
  (20) \qquad \psi_{\vec{a}} : \vec{\mathfrak{Z}}_s \xrightarrow{\ \sim\ }
  \vec{\mathfrak{Z}}_{s'}, \qquad (21) \qquad \psi_{\overleftarrow{a}} =
  (\psi_{\vec{a}})^{-1},
  \]
  qui dit comment l'orientation locale se transporte le long de l'arête.
\end{enumerate}
Ce sont les \emph{systèmes de rotation généralisés} (ou schémas de
plongement) de la théorie des plongements de graphes, qui codent les
plongements cellulaires dans les surfaces non nécessairement
orientables\footnote{Le nom est de nous, cité de mémoire (Ringel ; Stahl,
1978 ; Gross et Tucker, 1987). La page précise « quand les sommets de $K$
sont d'ordre $\geqslant 4$ (a fortiori $\geqslant 3$) » : à partir de $3$, les
deux ordres circulaires sont distincts. Le signe de (19) est lu $\subset$,
peut-être $\in$.}. « Canonique » veut dire : provenant de la carte. À
l'inverse, si l'on sait construire intrinsèquement, à partir de $(K, (D_i),
\omega)$, des données (19) et (20) satisfaisant (21), on obtient un
plongement cellulaire de $K$ dans une surface fermée $X$ --- « peut-être pas
le plan projectif ! ». Deux conditions s'imposent alors.

La première dit que les $D_i$ se déduisent de la carte de la façon
habituelle, en « allant tout droit » à chaque sommet :
\[
(22) \qquad \rho^{\nu_s}(\vec{D}_i) = -\vec{D}_i \quad \text{pour } \rho \in
\vec{\mathfrak{Z}}_s \text{ et } \vec{D}_i \text{ passant par } s :
\]
autour de $s$, l'arête opposée à une arête de $D_i$ est l'autre arête de
$D_i$, à mi-tour\footnote{En marge : si $\nu_s = 2$, (22) détermine déjà
$\vec{\mathfrak{Z}}_s$.}.

La seconde dit que la surface obtenue est le plan projectif. C'est la
relation d'Euler-Poincaré
\[
(23) \qquad c_0 - c_1 + c_2 = 1, \qquad \text{soit} \qquad c_2 = c(K),
\]
où $c_0$ et $c_1$, nombres de sommets et d'arêtes, ne dépendent que de $K$,
$c_2$ (nombre de faces) dépend de la carte, et $c(K) = c_1 - c_0 + 1$ est
le « degré de connexion » de $K$, son premier nombre de Betti. En marge :
on a a priori $c_2 \leqslant c(K)$. C'est juste, pour la raison suivante :
$K$ est connexe (deux pseudo-droites se rencontrent toujours), donc $X$
aussi ; et $X$ n'est pas une sphère, car deux courbes fermées qui se
coupent transversalement en un seul point ne peuvent pas être tracées sur une
sphère. Donc $\chi(X) \leqslant 1$, avec égalité si et seulement si $X$ est
un plan projectif\footnote{L'argument est de nous ; la note marginale ne se
lit qu'en partie.}.

\pagerange{42}{46}
\subsection*{L'ordre circulaire en un sommet (pages 42 à 46)}

Soit $s$ un sommet par lequel passent $\nu_s \geqslant 3$ pseudo-droites,
$I_s$ leur ensemble, et $\Delta$ une pseudo-droite qui ne passe pas par $s$.
L'application
\[
(25) \qquad \varphi_{s,\Delta} : I_s \longrightarrow S_\Delta, \qquad D_i \cap
\Delta = \{\varphi_{s,\Delta}(D_i)\},
\]
vers l'ensemble $S_\Delta$ des sommets de $K$ sur $\Delta$, est
\emph{injective} : deux pseudo-droites par $s$ ne peuvent se recouper sur
$\Delta$. Comme $S_\Delta$ est un polygone, $I_s$ reçoit une structure
polygonale induite $\mathfrak{Z}_{s,\Delta}$\footnote{La page note
$\mathfrak{Z}^{1}_s$ et $\mathfrak{Z}^{1}_{s,\Delta}$, l'exposant étant un petit
trait crochu lu « $1$ » sans certitude ; on l'omet, la flèche suffisant à
distinguer $\mathfrak{Z}_s$ de $\vec{\mathfrak{Z}}_s$. La première marge de la
page 42 l'écrit d'ailleurs sans exposant.}. On demande :
\[
(\mathrm{A}1) \qquad \mathfrak{Z}_{s,\Delta} \text{ ne dépend pas de } \Delta
\in I \smallsetminus I_s ; \text{ on la note } \mathfrak{Z}_s .
\]
C'est une condition sur les sommets d'au moins trois pseudo-droites ;
qu'elle soit satisfaite pour un $K$ provenant d'une carte, la marge le note
« à prouver ». Dans le plan, c'est l'énoncé familier que les droites d'un
faisceau coupent toute autre droite dans le même ordre.

Soit $\vec{I}_s \to I_s$ l'application de degré $2$ qui oublie l'orientation
(26). Si $\vec{u} \in \vec{\mathfrak{Z}}_s$, son image $u$ est une
permutation circulaire de $I_s$, et l'on doit avoir
\[
(27) \qquad u \in \mathfrak{Z}_s .
\]
C'est une condition restrictive, qui ne laisse que $2^{\nu_s - 1}$ relèvements
possibles d'un $u$ donné : relever $u$, c'est choisir pour chaque $i \in I_s$
une bijection $\omega_i \simeq \omega_{u(i)}$ entre les ensembles des deux
orientations de $D_i$ et de $D_{u(i)}$, soit $2^{\nu_s}$ choix, et l'on veut
que la permutation $\vec{u}$ obtenue de $\vec{I}_s$ soit circulaire,
c'est-à-dire que le composé des $\nu_s$ bijections autour du polygone soit
l'échange des deux orientations et non l'identité : les $\nu_s - 1$ premiers
choix fixent le dernier\footnote{Un premier essai de ce compte est encadré et
barré. La page désigne aussi ces ensembles d'orientations par $\omega_i$.}.

\textbf{La règle du triangle.} Soit maintenant $\vec{\Delta}$ une
pseudo-droite \emph{orientée} ne passant pas par $s$. Son orientation
oriente le polygone $S_\Delta$, donc $I_s$ : elle choisit un générateur
$u_{s,\vec{\Delta}}$ de $\mathfrak{Z}_s$ --- pour $\nu_s = 2$, c'est la
transposition de $I_s$, et $\mathfrak{Z}_s$ a alors un seul élément. Il
s'agit de relever $u_{s,\vec{\Delta}}$ en $\vec{u}_{s,\vec{\Delta}}$ : pour
$D_i$ passant par $s$ et $D_j$ son successeur, décider, une orientation
$\vec{D}_i$ étant donnée, l'orientation correspondante de $D_j$. La figure
(29) le dit : soient $t = D_i \cap \Delta$ et $u = D_j \cap \Delta$ ;
l'orientation de $D_j$ est celle pour laquelle l'arête $a$ de $s$ à $t$ dans
le sens de $\vec{D}_i$, l'arête $b$ de $s$ à $u$ dans le sens cherché de
$D_j$, et l'arête $c$ de $t$ à $u$ dans le sens de $\vec{\Delta}$ forment un
système d'orientations opposé à l'orientation $\omega$ du triple $\{D_i,
D_j, \Delta\}$. La règle ne fait intervenir que ces trois pseudo-droites et
$\omega$\footnote{La convention de signe, « opposé à l'orientation », est
celle de la page et dépend de la figure (29), que la transcription décrit
sans la reproduire ; on ne la recalcule pas. Le triangle est noté
$W_{\{i, j, \Delta\}}$, indice lu sans certitude.}. On demande :
\begin{itemize}
  \item[(A2)] la permutation $\vec{u}_{s,\vec{\Delta}}$ ainsi définie est
  circulaire, c'est-à-dire que $(\vec{u}_{s,\vec{\Delta}})^{\nu_s}$ est
  l'échange sur chaque fibre de (26), et non l'identité ;
  \item[(30)] on a alors $\vec{u}_{s,-\vec{\Delta}} =
  \vec{u}_{s,\vec{\Delta}}^{\,-1}$, et l'on pose $\vec{\mathfrak{Z}}_{s,\Delta}
  = \{\vec{u}_{s,\vec{\Delta}}, \vec{u}_{s,-\vec{\Delta}}\} \subset
  \mathcal{C}_{\vec{I}_s}$ ;
  \item[(A3)] pour toute autre pseudo-droite $\Delta'$ ne passant pas par
  $s$, $\vec{\mathfrak{Z}}_{s,\Delta'} = \vec{\mathfrak{Z}}_{s,\Delta}$ ; on le
  note $\vec{\mathfrak{Z}}_s$.
\end{itemize}
« Bien entendu », dit la page, (A1) à (A3) sont satisfaites quand
l'arrangement orienté provient d'un arrangement de dimension $2$ ; on
obtient alors les $\vec{\mathfrak{Z}}_s$, et une application surjective
\[
(31) \qquad \vec{I} \smallsetminus \vec{I}_s \longrightarrow
\vec{\mathfrak{Z}}_s, \qquad \vec{\Delta} \longmapsto
\vec{u}_{s,\vec{\Delta}},
\]
compatible à l'action de $\{\pm 1\}$.

\pagerange{47}{49}
\subsection*{Le transport le long d'une arête (pages 47 à 49)}

Soit $\vec{a}$ une arête de $K$ de $s$ à $s'$, et supposons qu'il existe une
pseudo-droite qui ne passe ni par $s$ ni par $s'$. Soit $\vec{J} = \vec{I}
\smallsetminus (\vec{I}_s \cup \vec{I}_{s'})$ ; les applications (31) en $s$
et en $s'$ font de $\vec{\mathfrak{Z}}_s$ et $\vec{\mathfrak{Z}}_{s'}$ deux
quotients de $\vec{J}$,

(32)

\begin{tikzcd}[column sep=small]
& \vec{J} \arrow[dl] \arrow[dr] & \\
\vec{\mathfrak{Z}}_s \arrow[rr, dashed, "\psi_{\vec{a}}"] & &
\vec{\mathfrak{Z}}_{s'}
\end{tikzcd}

et $\psi_{\vec{a}}$ existe si et seulement si ce sont les mêmes quotients :
\begin{itemize}
  \item[(A4)] si deux pseudo-droites orientées $\vec{\Delta}$, $\vec{\Delta}'$ ne
  passant ni par $s$ ni par $s'$ induisent la même orientation en $s$, elles
  induisent la même orientation en $s'$.
\end{itemize}

Reste le cas « très particulier » où toute pseudo-droite passe par $s$ ou par
$s'$. Soit $\Theta$ la pseudo-droite qui porte $a$, $D \neq \Theta$ une
pseudo-droite par $s$, $D' \neq \Theta$ une par $s'$, $t = D \cap D'$, et
$K' = K_{\{\Theta, D, D'\}}$, graphe octaédral gauche. Soit $b$ l'une des deux
arêtes de $K'$ joignant $s$ à $t$, et $b'$ l'unique arête de $K'$ joignant
$t$ à $s'$ telle que $\{a, b, b'\}$ soit compatible avec l'orientation de
$K'$, c'est-à-dire borde une face (figure (33)). Les arêtes $b$ et $b'$ de
$K'$ sont des chemins d'arêtes de $K$, portés par $D$ et $D'$. Excluons le
cas où $K$ est formé de $n \geqslant 3$ pseudo-droites concourantes et d'une
seule sécante : alors $\Theta$ est l'unique pseudo-droite non bigonale, les
transports le long des arêtes portées par $D$ et $D'$ sont déjà définis, et
l'on pose
\[
(34) \qquad \psi_{\vec{a}} = \psi_{\vec{b}'} \, \psi_{\vec{b}},
\]
où $\psi_{\vec{b}}$ et $\psi_{\vec{b}'}$ sont les composés des transports le
long des arêtes de $K$ qui composent les chemins de $s$ à $t$ et de $t$ à
$s'$\footnote{Les deux dernières lignes de la page 48 sont d'une lecture très
incertaine. Une note marginale ajoute que $K$ est alors réduit à un cas
« déjà traité ».}.

Il faudrait encore savoir que ce $\psi_{\vec{a}}$ ne dépend pas des choix
faits ; la page appelle cette condition (A5), nom qu'elle redonne page 51 à
une autre condition\footnote{On garde (A5) pour celle de la page 51, et on
ne numérote pas celle-ci.}. Et Grothendieck conclut :

\begin{quote}
« Le cas où $K$ comporte une ps.\ droite qui est un bigône devrait plutôt
être traité à part, avec des moyens plus directement adaptés à cette
situation. Je n'ai donc pas vérifié que la proposition énoncée soit correcte,
sans exclure le cas où $K$ est formé de $n$ pseudo-droites concourantes,
coupées par une seule autre ps.\ droite. »
\end{quote}

Le théorème de la page 39 reste donc, de son propre aveu, non démontré dans
ce cas. On ne sait pas s'il y est faux. On remarquera seulement que, dans ce
cas, l'automorphisme du squelette qui échange simultanément les deux arêtes
de chaque bigone respecte $\omega$ et provient bien de la carte : c'est
l'homologie harmonique de centre le point de concours et d'axe la sécante\footnote{Remarque de nous, qui ne règle pas la question.}.

\pagerange{49}{54}
\section*{VII. L'orientation canonique d'un triangle (pages 49 à 54)}

« J'ai maintenant envie de montrer que dans le cas d'un $K$ provenant d'une
$\mathcal{X}$, son orientation peut se décrire canoniquement en termes de
$(K, (D_i))$ seulement. » Si c'est le cas, la donnée $\omega$ du théorème de
la page 39 est superflue.

\pagerange{50}{51}
\subsection*{Premier cas : une quatrième pseudo-droite évite le triangle (pages 50 et 51)}

Soit $J \in \mathfrak{P}_3(I)$ un triple de pseudo-droites non
concourantes\footnote{La page écrit « concourantes » ; un triangle demande
« non concourantes », et l'on ne lit pas de « non » sur la page.}, et $K_J$
le « triangle » qu'elles forment. Supposons qu'il existe une pseudo-droite
$D_i$, $i \notin J$, qui ne passe par aucun des trois sommets de $K_J$ : les
quatre pseudo-droites de $J \cup \{i\}$ sont trois à trois non concourantes.
On utilise :

\emph{Proposition.} Le foncteur évident
\[
\begin{pmatrix} \text{arrangements simples} \\ \text{de } 4
\text{ pseudo-droites} \end{pmatrix} \longrightarrow
\begin{pmatrix} \text{arrangements de dimension } 1 \\ \text{simples de } 4
\text{ pseudo-droites} \end{pmatrix}
\]
est une équivalence de catégories --- on le voit en comparant les deux
groupoïdes à celui des ensembles à quatre éléments : chacun n'a qu'une
classe, de groupe $\mathfrak{S}_4$ agissant sur les pseudo-droites\footnote{La
justification par les groupes est de la page ; le détail --- un
automorphisme du squelette qui fixe les quatre pseudo-droites fixe les six
sommets, intersections de deux d'entre elles, donc tout --- est de nous.}.

Comme $K_{J \cup \{i\}}$ provient alors canoniquement d'un arrangement de
dimension $2$, $K_J$ en hérite une orientation $\omega_{J,i}$, et l'on demande
\begin{itemize}
  \item[(A5)] l'orientation $\omega_{J,i}$ de $K_J$ ne dépend pas du choix de $i$.
\end{itemize}

\pagerange{51}{54}
\subsection*{Second cas : toute pseudo-droite passe par un sommet du triangle (pages 51 à 54)}

Supposons que toute pseudo-droite de $K$ passe par l'un des trois sommets
$a_i$ ($i \in J$) de $K_J$, $a_i$ désignant le sommet opposé à $D_i$. Comme
il n'y a pas de bigone, par chaque sommet passe au moins une $D_{a_i}$ avec
$a_i \in D_{a_i}$, $D_{a_i} \notin J$.

Supposons d'abord qu'on puisse choisir $D_{a_1}, D_{a_2}, D_{a_3}$ non
concourantes. On obtient une configuration de six pseudo-droites, associée
canoniquement à l'ensemble $J$, à neuf sommets : les $a_i$, les points
$\alpha_i = D_{a_i} \cap D_i$, et les points $a'_i = D_{a_j} \cap D_{a_k}$ ($J
= \{i, j, k\}$). Chaque $D_{a_i}$ porte quatre sommets, $a_i$, $\alpha_i$,
$a'_j$, $a'_k$, et la configuration est entièrement décrite quand on sait
lequel des trois autres est \emph{opposé} à $a_i$ sur le quadrilatère
$D_{a_i}$. Deux cas sont possibles :
\begin{itemize}
  \item \emph{cas I} : il existe un ordre circulaire (unique) sur $J$ tel
  que $a_i$ soit opposé à $a'_{i+1}$ ;
  \item \emph{cas II} : $a_i$ est opposé à $\alpha_i$ pour tout $i$.
\end{itemize}
Soit $\lambda_i$ l'arête de $K_J$ portée par $D_i$ qui contient
$\alpha_i$. Si la configuration provient d'un arrangement de dimension $2$,
le triple $\{\lambda_i\}_{i \in J}$ est compatible avec l'orientation de
$K_J$ dans le cas I, et ne l'est pas dans le cas II\footnote{La page dit
d'abord que le cas d'une figure de dimension $2$ est le cas I, puis ajoute en
marge le cas II comme « autre possibilité ». Plusieurs mots du passage sont
illisibles ; la répartition des cas est celle que la transcription
établit.}. D'où la condition :
\begin{itemize}
  \item[(A6)] pour un triple $J$ non concourant tel que toute autre pseudo-droite
  passe par un sommet de $K_J$, et lorsqu'on peut choisir des $D_{a_i}$ non
  concourantes, on exige 1\textsuperscript{o}) que la configuration de six
  pseudo-droites obtenue soit du cas I ou du cas II, et 2\textsuperscript{o})
  que l'orientation de $K_J$ qu'elle définit ne dépende pas du choix des
  $D_{a_i}$.
\end{itemize}

Reste le cas où les $D_{a_i}$ sont concourantes quel que soit le choix.
L'arrangement est alors l'arrangement standard de six pseudo-droites joignant
deux à deux quatre points triples, avec $7 = 4 + 3$ sommets, les trois points
supplémentaires étant les intersections des pseudo-droites « opposées » : le
quadrangle complet $\mathrm{III}_6$ de la page 24\footnote{L'identification
avec $\mathrm{III}_6$ est notée par le transcripteur ; le dessin de la page
53 est un triangle et ses trois médianes, comme page 20.}. Sa forme de
dimension $1$ est décrite par l'ensemble de ses quatre points triples, et sa
forme de dimension $2$ aussi, de groupe d'automorphismes $\mathfrak{S}_4$. Si
$a, b, c$ sont les sommets de $K_J$, $\alpha, \beta, \gamma$ les sommets
« opposés » ($\alpha$ est l'intersection de $D_a$ avec la pseudo-droite qui
joint $b$ et $c$), et $\lambda_a, \lambda_b, \lambda_c$ les arêtes de $K_J$ qui
contiennent $\alpha, \beta, \gamma$, alors $\{\lambda_a, \lambda_b,
\lambda_c\}$ est compatible avec l'orientation.

\textbf{En résumé.} \emph{Théorème} (page 54). \emph{Le foncteur canonique du
groupoïde des arrangements de pseudo-droites sans bigone vers celui des
arrangements de dimension $1$ de pseudo-droites sans bigone est pleinement
fidèle.}

Il s'obtient en composant le théorème de la page 39 avec l'orientation
canonique qu'on vient de décrire : tout isomorphisme de squelettes respecte
les orientations canoniques, puisqu'elles sont définies à partir du seul
squelette. Le cas litigieux de la page 49 a des bigones et n'intervient
pas ; restent non vérifiés l'indépendance des choix dans (34), et les
conditions (A1) à (A6) pour un arrangement de dimension $2$, affirmées sans
démonstration. Grothendieck ajoute que (A1) à (A6) donnent en principe une
voie pour déterminer l'image essentielle, mais qu'il faudrait d'abord les
simplifier, et qu'il faudrait encore y joindre la relation d'Euler-Poincaré
(23), « mais je soupçonne que cette dernière est automatique ».

\pagerange{56}{58}
\section*{Dessins (pages 56 à 58)}

Trois petites cartes à l'italienne. La page 56 dessine un arrangement de
droites à cinq sommets multiples, chaque région marquée de son nombre de
côtés, les régions à l'infini marquées « id », et compte $25$ triangles, $5$
carrés et $1$ pentagone, soit $31$ régions. La page 57, six droites qui
forment au centre un petit triangle entouré de triangles. La page 58, un
arrangement d'une dizaine de droites dont les extrémités portent des
étiquettes de types, et le décompte
\[
5 + (5 + 10) + (5 + 5) + 1 = 31,
\]
avec la légende : $1_6$, « icosaèdre » ; $5_1$, « hexagone » ; $5_4$ et
$5_2$, « triangle bordé » ; $5_5$ et $10_3$, « carré
gammé »\footnote{« gammé » est une lecture incertaine. Aucun arrangement
simple n'a $31$ régions ($1 + N(N-1)/2$ ne prend pas cette valeur) : les
arrangements de ces pages ont des points multiples. Remarque de nous ; on ne
tente pas d'identifier les figures.}.

\pagerange{60}{66}
\section*{VIII. « Système de Pseudodroites » : un second départ (pages 60 à 66)}

\pagerange{60}{61}
\subsection*{Définition et premières conséquences (pages 60 et 61)}

\emph{Définition.} Un \emph{système de pseudo-droites} est une carte,
topologique ou combinatoire, dont la surface sous-jacente $X$ est un plan
projectif réel (une surface compacte non orientable de caractéristique
d'Euler-Poincaré $1$), dont tous les sommets sont d'ordre pair $\geqslant 4$
--- ce qui partage l'ensemble des arêtes en classes d'équivalence, en allant
tout droit à chaque sommet : ce sont les \emph{pseudo-droites} de la carte
---, et telle que deux pseudo-droites se rencontrent en un point et un
seul\footnote{Il faut aussi demander que chaque pseudo-droite passe au plus
une fois par chaque sommet, c'est-à-dire soit une courbe fermée simple ;
« deux pseudo-droites se rencontrent en un point et un seul » ne l'entraîne
pas. L'hypothèse est sous-entendue sur la page.}.

\emph{Premières conséquences.} Le nombre $n$ des pseudo-droites est
$\geqslant 2$. Chaque pseudo-droite est une courbe non homotope à $0$ : le
long d'elle, $X$ est non orientable, et la situation est la même que pour
une droite du plan projectif\footnote{Cela découle de ce que deux
pseudo-droites se coupent transversalement en un seul point : leur nombre
d'intersection modulo $2$ vaut $1$, ce qui interdit à l'une ou l'autre
d'être homologue à zéro. L'argument est de nous.}. « On montre » que pour
$n \leqslant 6$ pseudo-droites en position générale (tous les sommets d'ordre
$4$), le système peut être réalisé par de vraies droites. C'est vrai, et
même mieux : tout arrangement d'au plus huit pseudo-droites est
\emph{étirable} (Goodman et Pollack, 1980), et il existe des arrangements
simples non étirables de neuf pseudo-droites (Ringel, 1956)\footnote{Les
références sont citées de mémoire ; la page n'en cite aucune et ne donne pas
la démonstration.}.

Soit $D$ une pseudo-droite, $\widetilde{X}$ le revêtement des orientations de
$X$ --- une sphère, $\widetilde{X} \simeq \mathbf{S}^2 \simeq
\mathbf{P}^1_{\mathbf{C}}$ --- et $\widetilde{D}$ l'image inverse de $D$ dans
$\widetilde{X}$\footnote{La page écrit « l'image inverse de $D$ dans $X$ ».},
un cercle. Alors $\widetilde{X} \smallsetminus \widetilde{D}$ a deux composantes
connexes, deux « hémisphères », en bijection canonique avec les deux
orientations de $D$ : un hémisphère est une orientation du disque $X
\smallsetminus D$, qui oriente son bord, revêtement double de $D$, donc
$D$\footnote{L'explication de la bijection est de nous ; la page l'affirme.}.

Les sommets portés par $D$ forment un polygone combinatoire $S_D$ dès qu'ils
sont au moins trois, et les orientations de $D$ --- c'est-à-dire
$\pi_0(\widetilde{X} \smallsetminus \widetilde{D})$ --- correspondent à celles de
ce polygone. La marge propose de dire plutôt que $D$ est le support d'un
sous-complexe polygonal du squelette $K^1$, sans exclure les monogones et les
bigones. Notons $\Omega(D)$ l'ensemble des deux orientations de $D$.

\pagerange{61}{64}
\subsection*{Comparer les orientations de deux pseudo-droites (pages 61 à 64)}

Soit $s$ un sommet. Si toutes les pseudo-droites ne passent pas par $s$, il
en existe une, $D$, qui n'y passe pas ; $s$ est alors dans le disque $X
\smallsetminus D$, et l'ensemble $\Omega_s(X)$ des orientations locales de $X$
en $s$ est en bijection canonique avec $\Omega(D)$. Si $D'$ est une autre
pseudo-droite ne passant pas par $s$, on a donc une bijection composée
\[
\Omega(D) \simeq \Omega_s(X) \simeq \Omega(D'),
\]
et la question est de l'expliciter en termes du seul squelette $K^1$ et de
ses pseudo-droites.

Soit $t = D \cap D'$. Comme $s$ est d'ordre $\geqslant 4$, il passe par $s$
au moins deux pseudo-droites $\Delta, \Delta'$ ; supposons-les distinctes de
$D, D'$ et ne passant pas par $t$, et soient $u, u'$ (resp.\ $v, v'$) leurs
intersections avec $D$ et $D'$. Les quatre pseudo-droites $D, D', \Delta,
\Delta'$ sont en position générale, donc en position standard, et l'on
constate que les orientations « de $t$ vers $u$ » de $D$ et « de $t$ vers
$u'$ » de $D'$ se correspondent. Plus généralement, si $\Delta_1, \ldots,
\Delta_k$ passent par $s$ et non par $t$, coupant $D$ en $u_i$ et $D'$ en
$u'_i$, et si l'on munit $D \smallsetminus \{t\}$ et $D' \smallsetminus \{t\}$
des ordres définis par deux orientations qui se correspondent, alors les
$u'_i$ croissent quand les $u_i$ croissent : c'est la perspective de centre
$s$. Cela définit la bijection $\Omega(D) \simeq \Omega(D')$
combinatoirement.

Reste le cas où une seule pseudo-droite $\Delta$ passe par $s$ et non par
$t$ : alors $s$ est d'ordre $4$, et l'autre pseudo-droite $\Delta'$ par $s$
passe par $t$. S'il existe une pseudo-droite $D''$ qui ne passe pas par $t$,
distincte de $\Delta$, coupant $D$ et $D'$ en $x \neq u$ et $x' \neq u'$, on
compare $\Omega(D)$ et $\Omega(D'')$, puis $\Omega(D'')$ et $\Omega(D')$,
« et on gagne ». Si toutes les pseudo-droites qui ne passent pas par $t$
passent par $u$, ou toutes par $u'$, on a deux faisceaux de centres $t$ et
$u$ (resp.\ $u'$), cas qui demande un traitement à part. Sinon, il existe
$\Gamma$ passant par $u$ et non par $u'$ et $\Gamma'$ passant par $u'$ et non
par $u$ ; on a $\Omega(\Gamma) \simeq \Omega(D')$ et $\Omega(\Gamma') \simeq
\Omega(D)$, et l'on peut déduire $\Omega(D) \simeq \Omega(D')$ de
$\Omega(\Gamma) \simeq \Omega(\Gamma')$ pourvu que l'intersection $r$ de
$\Gamma$ et $\Gamma'$ ne soit pas sur $\Delta'$\footnote{La fin de la
condition n'est pas écrite : un blanc suit « l'intersection $r$ de $\Gamma$
et $\Gamma'$ ». On la complète d'après la phrase suivante, « par malheur on
se trouve encore justement sur $\Delta'$ ».}. Si « par malheur » $r$ est sur
$\Delta'$, les six pseudo-droites $D, D', \Delta, \Delta', \Gamma, \Gamma'$
sont celles qui joignent deux à deux les quatre points $u, u', r, t$ --- de
nouveau le quadrangle complet --- et l'on constate que la bijection
$\varphi_s : \Omega(D) \simeq \Omega(D')$ définie par $s$ est celle que
définit $r$, qui est dans la même des deux régions de $X$ découpées par $D$
et $D'$.

\pagerange{64}{66}
\subsection*{Axiomes et système local (pages 64 à 66)}

Soit $K^1$ un $1$-complexe découpé en pseudo-droites --- tout sommet sur au
moins deux d'entre elles, deux pseudo-droites se rencontrant en un point et
un seul. Pour deux pseudo-droites distinctes $D, D'$ et un sommet $s$ qui
n'est sur aucune des deux, on se donne une bijection
\[
\varphi_{s;D,D'} : \Omega(D) \simeq \Omega(D'),
\]
avec les propriétés suivantes.
\begin{itemize}
  \item[a)] Sur l'ensemble des pseudo-droites qui ne passent pas par $s$, les
  $\varphi_{s;D,D'}$ forment un système transitif d'isomorphismes :
  $\varphi_{s;D',D''} \circ \varphi_{s;D,D'} = \varphi_{s;D,D''}$\footnote{La
  page écrit « est simplement transitif », où « simplement » est une lecture
  incertaine ; ce qui est demandé, et ce que la suite utilise pour passer au
  quotient, est la compatibilité à la composition.}.
  \item[b)] Si $\Delta$ est une pseudo-droite distincte de $D$ et $D'$, les
  coupant en $x$ et $x'$ (on n'exclut pas $x = x'$), et si $s, s'$ sont deux
  sommets de $\Delta$ distincts de $x, x'$, alors $\varphi_{s;D,D'} =
  \varphi_{s';D,D'}$ si et seulement si $s$ et $s'$ sont sur un même des deux
  arcs que $x$ et $x'$ découpent sur $\Delta$. On peut donc définir
  $\varphi_{a;D,D'}$ pour un tel arc $a$.
  \item[c)] Soit $t = D \cap D'$, et $\Delta_1, \Delta_2$ deux pseudo-droites par $s$
  et non par $t$, coupant $D$ en $u_1, u_2$ et $D'$ en $u'_1, u'_2$ \ldots{} ---
  la page s'arrête là.
\end{itemize}
Dans un arrangement de dimension $2$, b) dit que $\varphi_{s;D,D'}$ ne dépend
que de celle des deux régions de $X \smallsetminus (D \cup D')$ qui contient
$s$ ; quand $x = x'$, l'arc unique $\Delta \smallsetminus \{t\}$ est tout
entier dans une région\footnote{Cette lecture géométrique de b) est de
nous.}.

La page 65 commence sur une phrase dont le début manque : le système
d'isomorphismes en question serait unique, à certaines conditions. Il ne
peut pas l'être s'il existe une pseudo-droite bigonale --- c'est-à-dire un
double faisceau avec la pseudo-droite qui joint ses deux centres --- ni
s'il existe un monogone, c'est-à-dire un seul faisceau.

Dans le cas où il n'y a ni monogone ni bigone, on définit pour tout sommet
$s$ l'ensemble $\Omega(s)$ comme le quotient de la réunion des $\Omega(D)$,
$D$ ne passant pas par $s$, par le système des $\varphi_{s;D,D'}$ : c'est la
limite d'un système transitif d'isomorphismes, un ensemble à deux éléments.
Pour en faire un \emph{système local} sur $K^1$, il faut pour toute arête de
$K^1$, d'extrémités $s$ et $s'$, une bijection $\Omega(s) \simeq \Omega(s')$ ;
on prend une pseudo-droite $D$ qui ne passe ni par $s$ ni par $s'$, et la
composée

\begin{tikzcd}[column sep=small]
\Omega(s) \arrow[rr, "\sim"] \arrow[dr, "\sim"'] & & \Omega(s') \\
& \Omega(D) \arrow[ur, "\sim"'] &
\end{tikzcd}

qui ne dépend pas du choix de $D$ par b) --- c'en est même une forme
équivalente\footnote{La page dit qu'une telle $D$ existe « par hyp. » ; c'est
une hypothèse supplémentaire : le cas où toute pseudo-droite passe par l'une
des extrémités de l'arête est celui que la page 48 traitait à part.}.

Pour reconstituer la carte à partir de $K^1$, il faudrait encore définir une
structure polygonale sur l'ensemble des arcs issus de chaque sommet $s$, et
une bijection canonique entre $\Omega(s)$ et les deux orientations de cette
structure. La page 66 s'arrête là.

\pagerange{67}{70}
\section*{IX. Pseudo-droites orientées : un troisième départ (pages 67 à 70)}

Sur de plus petites feuilles, la même construction est reprise avec les
pseudo-droites orientées. Soit $\vec{\mathcal{D}}$ l'ensemble des
pseudo-droites orientées, et $\vec{\mathcal{D}}(s)$ celui des pseudo-droites
orientées passant par $s$ --- autant que d'arcs issus de $s$, soit $2\nu(s)$.
Pour $\vec{D}$ ne passant pas par $s$, on se donne une permutation circulaire
$\rho_{\vec{D},s}$ de $\vec{\mathcal{D}}(s)$\footnote{La lettre, un long trait
bouclé en haut, est lue $\rho$ sans certitude ; de même le $\nu$ de
« cardinal $2\nu(s)$ ».}, avec les axiomes :
\begin{itemize}
  \item[a)] $\rho_{-\vec{D},s} = \rho_{\vec{D},s}^{-1}$ ;
  \item[b)] si $\vec{D}, \vec{D}' \notin \vec{\mathcal{D}}(s)$, $\rho_{\vec{D},s}$ et
  $\rho_{\vec{D}',s}$ sont égales ou inverses l'une de l'autre\footnote{« ou
  opposées », lu sans certitude.} ;
  \item[c)] soient $\vec{D}, \vec{D}', \vec{\Delta}$ des pseudo-droites
  distinctes, $x, x'$ les intersections de $\Delta$ avec $D$ et $D'$ (on ne
  suppose pas $x \neq x'$), $s, s'$ des sommets de $\Delta$ distincts de $x,
  x'$, et supposons $\rho_{\vec{D},s} = \rho_{\vec{D}',s}$ ; pour qu'on ait
  $\rho_{\vec{D},s'} = \rho_{\vec{D}',s'}$, il faut et il suffit que $s$ et
  $s'$ soient sur un même arc de $\Delta$ découpé par $x$ et
  $x'$\footnote{Un premier état de c), page 67, est barré et repris page
  68.} ;
  \item[(d)] soient $\vec{\Delta}_1 \in \vec{\mathcal{D}}(s)$, $\vec{D} \notin
  \vec{\mathcal{D}}(s)$, $\vec{\Delta}_2 = \rho_{\vec{D},s}(\vec{\Delta}_1)$,
  et $u_1$, $u_2$ les points où $\Delta_1$, $\Delta_2$ coupent $D$ : alors
  $u_2$ est le sommet de $\vec{D}$ qui suit $u_1$ \emph{parmi} les
  intersections de $\vec{D}$ avec les pseudo-droites passant par $s$ ;
  \item[(e)] soient $\vec{D}, \vec{D}' \notin \vec{\mathcal{D}}(s)$ distinctes,
  $t = D \cap D'$, $\Delta_1, \Delta_2$ deux pseudo-droites par $s$ et non par
  $t$, coupant $D$ en $u_1, u_2$ et $D'$ en $u'_1, u'_2$. Si $u_1$ et $u_2$ sont
  adjacents sur $D$, $u'_1$ et $u'_2$ le sont sur $D'$. Si $\vec{D}$ est
  orientée dans le sens $t u_1 u_2$, $\vec{D}'$ dans le sens $t u'_1 u'_2$, et
  si $u_1, u_2$ sont adjacents parmi les intersections de $D$ avec les
  pseudo-droites par $s$, alors $\rho_{\vec{D},s}$ envoie la pseudo-droite
  orientée $s u_1 u'_1$ sur $s u_2 u'_2$\footnote{La page écrit « qui passent
  par $\vec{D}$ » où il faut « par $s$ », comme en (d), et note l'indice de
  $\rho$ sans le $s$.} ;
  \item[(f)] si $s$ est d'ordre $2n$, $(\rho_{\vec{D},s})^n(\vec{\Delta}) =
  -\vec{\Delta}$.
\end{itemize}
Ces axiomes sont ceux des pages 42 à 47 réécrits pour les
$\rho_{\vec{D},s}$ : $\rho_{\vec{D},s}$ est la permutation
$\vec{u}_{s,\vec{D}}$ de la page 45, (d) est la façon dont elle a été
définie, a) est (30), b) est (A3), c) est la propriété b) de la page 64, et (f) est (22)\footnote{Ces correspondances sont
de nous.}.

\emph{Th} (page 70, énoncé). \emph{S'il existe un système de
$\rho_{\vec{D},s}$ satisfaisant a) à f), et s'il n'y a pas de pseudo-droite
qui soit un monogone ou un bigone, ce système est unique. De plus, les
$\Omega(s) = \{\rho_s, \rho_s^{-1}\}$ forment de façon unique un système
local} --- la fin de la phrase est illisible. En marge : il suffit de d), e)
et f), et sans doute a) en est une conséquence ; pour donner un sens au
système local, il faut des hypothèses plus fortes que b) et c). L'unicité
est plausible sur l'axiome (d) seul pour ce qui est des pseudo-droites sous-jacentes :
celles qui passent par $s$ coupent $D$ en des points distincts, et (d) dit
dans quel ordre on les parcourt ; ce sont (e) et (f) qui fixent les
orientations\footnote{Remarque de nous, qui ne démontre pas le « Th ».}.

\pagerange{72}{72}
\section*{X. Petits arrangements : la géométrie et la combinatoire (page 72)}

La dernière page est un tableau. Pour chaque petit arrangement de droites du
plan projectif réel, il note s'il est « standard » au sens projectif ---
toutes ses réalisations sont projectivement équivalentes --- ou seulement au
sens combinatoire --- il n'a qu'une classe combinatoire ---, et compare le
groupe des transformations projectives qui le stabilisent, modulo sa
composante neutre, au groupe de ses automorphismes combinatoires.

\begin{itemize}
  \item \emph{Une droite, deux droites} : standard projectif.
  \item \emph{Trois droites en position générale} : standard projectif ; le groupe
  projectif modulo sa composante neutre, le groupe combinatoire et le groupe
  des permutations des quatre faces sont isomorphes --- c'est
  $\mathfrak{S}_4$, le groupe de l'hémi-octaèdre des pages 31 à 38.
  \item \emph{Trois droites concourantes} : standard projectif ; le groupe projectif
  modulo sa composante neutre est une extension de $\mathfrak{S}_3$ par
  $\mathbf{Z}/2$, isomorphe au groupe combinatoire, extension de $\mathbf{Z}/2$ par
  $\mathbf{Z}/6$ --- le groupe de l'hexagone, celui des six demi-droites autour
  du point de concours.
  \item \emph{Quatre droites en position générale} : standard projectif ; groupe
  combinatoire $\mathfrak{S}_4$\footnote{La page écrit « endom proj $\simeq$
  Id mais endom comb.\ $\simeq \mathfrak{S}_4$ ». Le groupe des transformations
  projectives qui stabilisent l'\emph{ensemble} des quatre droites est
  $\mathfrak{S}_4$ lui aussi, toute permutation de quatre droites en position
  générale étant réalisée par une unique transformation projective ; c'est
  le groupe de celles qui fixent \emph{chaque} droite qui est trivial.}.
  \item \emph{Quatre droites dont trois concourantes} : standard projectif ;
  extension de $\mathfrak{S}_3$ par $\mathbf{Z}/2$, isomorphe au groupe
  combinatoire, groupe de l'hexagone.
  \item \emph{Quatre droites concourantes} : pas standard projectif (le birapport
  des quatre droites est un module), mais standard combinatoire ; groupe
  combinatoire, celui de l'octogone (?).
  \item \emph{Cinq droites en position générale} : pas standard projectif, mais
  standard combinatoire ; groupe combinatoire, celui du
  pentagone\footnote{La page biffe le « pas » de « pas standard proj., mais
  standard comb. ». Il est nécessaire : cinq droites en position générale
  dépendent de $2 \cdot 5 - 8 = 2$ paramètres projectifs, et le « mais » qui
  suit le demande. Le groupe du pentagone, diédral d'ordre $10$, est celui de
  la page 20, « $(D_5)$ ».}.
  \item \emph{Six droites ou plus en position générale} : \emph{pas} standard
  combinatoire --- il y a plusieurs arrangements simples de six droites.
\end{itemize}
Le tableau mesure l'écart entre les deux points de vue du dossier : en
position générale, la combinatoire détermine l'arrangement jusqu'à cinq
droites et cesse de le faire à six ; la géométrie projective, elle, a des
modules dès quatre droites concourantes et dès cinq droites en position
générale.

\end{document}
